% paper03-descent-obstructions.tex
% Sheaf Descent and Čech Obstructions for Program Proofs
\documentclass[11pt]{article}
% jugeo-common.tex — shared preamble for all Judgment Geometry papers
% Include via: % jugeo-common.tex — shared preamble for all Judgment Geometry papers
% Include via: % jugeo-common.tex — shared preamble for all Judgment Geometry papers
% Include via: \input{jugeo-common}

% ── Packages ────────────────────────────────────────────────────────────
\usepackage{amsmath,amssymb,amsthm,mathtools}
\usepackage{tikz-cd}
\usepackage{listings}
\usepackage{xcolor}
\usepackage{xspace}
\usepackage{booktabs}
\usepackage{multirow}
\usepackage{enumitem}
\usepackage[expansion=false]{microtype}
\usepackage{hyperref}
\usepackage{cleveref}
% \usepackage{thmtools}  % disabled: not available in this TeX installation
\usepackage{float}
\usepackage{caption}
\usepackage{subcaption}
\usepackage{tabularx}
\usepackage{stmaryrd}       % \llbracket, \rrbracket
\usepackage{bbm}            % \mathbbm{1}

% ── Page geometry ───────────────────────────────────────────────────────
\usepackage[margin=1in]{geometry}

% ── Theorem environments ────────────────────────────────────────────────
\theoremstyle{plain}
\newtheorem{theorem}{Theorem}[section]
\newtheorem{lemma}[theorem]{Lemma}
\newtheorem{proposition}[theorem]{Proposition}
\newtheorem{corollary}[theorem]{Corollary}

\theoremstyle{definition}
\newtheorem{definition}[theorem]{Definition}
\newtheorem{example}[theorem]{Example}
\newtheorem{construction}[theorem]{Construction}

\theoremstyle{remark}
\newtheorem{remark}[theorem]{Remark}
\newtheorem{notation}[theorem]{Notation}

% ── Python listings style ──────────────────────────────────────────────
\definecolor{jugeoBlue}{HTML}{2563EB}
\definecolor{jugeoGreen}{HTML}{059669}
\definecolor{jugeoRed}{HTML}{DC2626}
\definecolor{jugeoPurple}{HTML}{7C3AED}
\definecolor{jugeoGray}{HTML}{6B7280}
\definecolor{jugeoBg}{HTML}{F8FAFC}

\lstdefinestyle{jugeo-python}{
  language=Python,
  basicstyle=\ttfamily\small,
  keywordstyle=\color{jugeoBlue}\bfseries,
  stringstyle=\color{jugeoGreen},
  commentstyle=\color{jugeoGray}\itshape,
  numberstyle=\tiny\color{jugeoGray},
  backgroundcolor=\color{jugeoBg},
  numbers=left,
  numbersep=6pt,
  frame=single,
  framerule=0.4pt,
  rulecolor=\color{jugeoGray!40},
  breaklines=true,
  breakatwhitespace=true,
  showstringspaces=false,
  tabsize=4,
  xleftmargin=1.5em,
  framexleftmargin=1.5em,
  morekeywords={dataclass,frozen,field,Enum,IntEnum,auto,Any,
                Mapping,Sequence,Optional,match,case},
  literate={→}{{$\to$}}1 {←}{{$\leftarrow$}}1
           {≤}{{$\leq$}}1 {≥}{{$\geq$}}1 {≠}{{$\neq$}}1
           {∈}{{$\in$}}1 {∀}{{$\forall$}}1 {∃}{{$\exists$}}1
           {⊕}{{$\oplus$}}1 {⊖}{{$\ominus$}}1,
  escapeinside={(*@}{@*)},
}
\lstset{style=jugeo-python}

\lstdefinestyle{jugeo-output}{
  basicstyle=\ttfamily\small,
  backgroundcolor=\color{jugeoBg},
  frame=single,
  framerule=0.4pt,
  rulecolor=\color{jugeoGray!40},
  breaklines=true,
  showstringspaces=false,
  xleftmargin=1.5em,
  framexleftmargin=0.5em,
  numbers=none,
}

% ── JuGeo-specific macros ──────────────────────────────────────────────

% Judgment 8-tuple
\newcommand{\J}{\mathcal{J}}
\newcommand{\Jud}[8]{(#1,\,#2,\,#3,\,#4,\,#5,\,#6,\,#7,\,#8)}
\newcommand{\JudShort}[1]{\J_{#1}}

% Components
\newcommand{\coord}{\mathsf{c}}            % coordinate
\newcommand{\prop}{\varphi}                 % proposition
\newcommand{\carrier}{\mathcal{A}}          % carrier type
\newcommand{\evid}{\mathcal{E}}             % evidence bundle
\newcommand{\oblig}{\mathcal{O}}            % obligations
\newcommand{\obstr}{\mathcal{B}}            % obstructions
\newcommand{\trust}{\mathcal{T}}            % trust annotation
\newcommand{\prov}{\Pi}                     % provenance

% Trust algebra
\newcommand{\Talg}{\mathfrak{T}}            % trust ordered algebra
\newcommand{\Eadm}{\mathcal{E}_{\mathrm{adm}}}
\newcommand{\tleq}{\preceq}                 % trust ordering
\newcommand{\tjoin}{\oplus}                 % conservative join
\newcommand{\tatten}{\ominus}               % attenuation
\newcommand{\tpromote}{\uparrow_\pi}        % promotion
\newcommand{\tdemote}{\downarrow_\chi}       % demotion

% Trust tiers
\newcommand{\Tcontra}{\ensuremath{\mathsf{CONTRADICTED}}}
\newcommand{\Tunver}{\ensuremath{\mathsf{UNVERIFIED}}}
\newcommand{\Tcopilot}{\ensuremath{\mathsf{COPILOT}}}
\newcommand{\Toracle}{\ensuremath{\mathsf{ORACLE}}}
\newcommand{\Truntime}{\ensuremath{\mathsf{RUNTIME}}}
\newcommand{\Tsolver}{\ensuremath{\mathsf{SOLVER}}}
\newcommand{\Tproof}{\ensuremath{\mathsf{PROOF}}}

% Site and geometry
\newcommand{\Site}{\mathcal{S}}             % semantic site
\newcommand{\Cov}{\mathrm{Cov}}            % covering family
\newcommand{\Ob}{\mathrm{Ob}}              % objects
\newcommand{\Mor}{\mathrm{Mor}}            % morphisms
\newcommand{\res}{\mathrm{res}}            % restriction
\newcommand{\Cech}{\check{C}}              % Čech complex
\newcommand{\Hcoh}[1]{H^{#1}}             % cohomology group

% Descent
\newcommand{\descent}{\mathrm{desc}}
\newcommand{\glue}{\mathrm{glue}}
\newcommand{\overlap}[2]{#1 \cap #2}

% Semantic moves
\newcommand{\VERIFY}{\textsc{Verify}}
\newcommand{\CONSTRUCT}{\textsc{Construct}}
\newcommand{\NORMALIZE}{\textsc{Normalize}}
\newcommand{\GLUE}{\textsc{Glue}}
\newcommand{\RESTRICT}{\textsc{Restrict}}
\newcommand{\TRANSPORT}{\textsc{Transport}}
\newcommand{\REPAIR}{\textsc{Repair}}
\newcommand{\PROMOTE}{\textsc{Promote}}
\newcommand{\CHALLENGE}{\textsc{Challenge}}
\newcommand{\ESCALATE}{\textsc{Escalate}}

% Fragments
\newcommand{\QFLIA}{\texttt{QF\_LIA}}
\newcommand{\QFLRA}{\texttt{QF\_LRA}}
\newcommand{\QFBV}{\texttt{QF\_BV}}
\newcommand{\QFUF}{\texttt{QF\_UF}}

% Misc
\newcommand{\Python}{\textsc{Python}}
\newcommand{\JuGeo}{\textsc{JuGeo}}
\newcommand{\LEAN}{\textsc{Lean}\,4}
\newcommand{\Fstar}{F$^{\star}$}
\newcommand{\Coq}{\textsc{Coq}}
\newcommand{\Dafny}{\textsc{Dafny}}

% Common shortcuts
\newcommand{\ie}{i.e.\@\xspace}
\newcommand{\eg}{e.g.\@\xspace}
\newcommand{\cf}{cf.\@\xspace}
\newcommand{\wrt}{w.r.t.\@\xspace}

% Evaluation shortcuts
\newcommand{\numcases}{300}
\newcommand{\numspec}{100}
\newcommand{\numeq}{100}
\newcommand{\numbug}{100}
\newcommand{\medtime}{6.5\,ms}
\newcommand{\totaltime}{1.949\,s}


% ── Packages ────────────────────────────────────────────────────────────
\usepackage{amsmath,amssymb,amsthm,mathtools}
\usepackage{tikz-cd}
\usepackage{listings}
\usepackage{xcolor}
\usepackage{xspace}
\usepackage{booktabs}
\usepackage{multirow}
\usepackage{enumitem}
\usepackage[expansion=false]{microtype}
\usepackage{hyperref}
\usepackage{cleveref}
% \usepackage{thmtools}  % disabled: not available in this TeX installation
\usepackage{float}
\usepackage{caption}
\usepackage{subcaption}
\usepackage{tabularx}
\usepackage{stmaryrd}       % \llbracket, \rrbracket
\usepackage{bbm}            % \mathbbm{1}

% ── Page geometry ───────────────────────────────────────────────────────
\usepackage[margin=1in]{geometry}

% ── Theorem environments ────────────────────────────────────────────────
\theoremstyle{plain}
\newtheorem{theorem}{Theorem}[section]
\newtheorem{lemma}[theorem]{Lemma}
\newtheorem{proposition}[theorem]{Proposition}
\newtheorem{corollary}[theorem]{Corollary}

\theoremstyle{definition}
\newtheorem{definition}[theorem]{Definition}
\newtheorem{example}[theorem]{Example}
\newtheorem{construction}[theorem]{Construction}

\theoremstyle{remark}
\newtheorem{remark}[theorem]{Remark}
\newtheorem{notation}[theorem]{Notation}

% ── Python listings style ──────────────────────────────────────────────
\definecolor{jugeoBlue}{HTML}{2563EB}
\definecolor{jugeoGreen}{HTML}{059669}
\definecolor{jugeoRed}{HTML}{DC2626}
\definecolor{jugeoPurple}{HTML}{7C3AED}
\definecolor{jugeoGray}{HTML}{6B7280}
\definecolor{jugeoBg}{HTML}{F8FAFC}

\lstdefinestyle{jugeo-python}{
  language=Python,
  basicstyle=\ttfamily\small,
  keywordstyle=\color{jugeoBlue}\bfseries,
  stringstyle=\color{jugeoGreen},
  commentstyle=\color{jugeoGray}\itshape,
  numberstyle=\tiny\color{jugeoGray},
  backgroundcolor=\color{jugeoBg},
  numbers=left,
  numbersep=6pt,
  frame=single,
  framerule=0.4pt,
  rulecolor=\color{jugeoGray!40},
  breaklines=true,
  breakatwhitespace=true,
  showstringspaces=false,
  tabsize=4,
  xleftmargin=1.5em,
  framexleftmargin=1.5em,
  morekeywords={dataclass,frozen,field,Enum,IntEnum,auto,Any,
                Mapping,Sequence,Optional,match,case},
  literate={→}{{$\to$}}1 {←}{{$\leftarrow$}}1
           {≤}{{$\leq$}}1 {≥}{{$\geq$}}1 {≠}{{$\neq$}}1
           {∈}{{$\in$}}1 {∀}{{$\forall$}}1 {∃}{{$\exists$}}1
           {⊕}{{$\oplus$}}1 {⊖}{{$\ominus$}}1,
  escapeinside={(*@}{@*)},
}
\lstset{style=jugeo-python}

\lstdefinestyle{jugeo-output}{
  basicstyle=\ttfamily\small,
  backgroundcolor=\color{jugeoBg},
  frame=single,
  framerule=0.4pt,
  rulecolor=\color{jugeoGray!40},
  breaklines=true,
  showstringspaces=false,
  xleftmargin=1.5em,
  framexleftmargin=0.5em,
  numbers=none,
}

% ── JuGeo-specific macros ──────────────────────────────────────────────

% Judgment 8-tuple
\newcommand{\J}{\mathcal{J}}
\newcommand{\Jud}[8]{(#1,\,#2,\,#3,\,#4,\,#5,\,#6,\,#7,\,#8)}
\newcommand{\JudShort}[1]{\J_{#1}}

% Components
\newcommand{\coord}{\mathsf{c}}            % coordinate
\newcommand{\prop}{\varphi}                 % proposition
\newcommand{\carrier}{\mathcal{A}}          % carrier type
\newcommand{\evid}{\mathcal{E}}             % evidence bundle
\newcommand{\oblig}{\mathcal{O}}            % obligations
\newcommand{\obstr}{\mathcal{B}}            % obstructions
\newcommand{\trust}{\mathcal{T}}            % trust annotation
\newcommand{\prov}{\Pi}                     % provenance

% Trust algebra
\newcommand{\Talg}{\mathfrak{T}}            % trust ordered algebra
\newcommand{\Eadm}{\mathcal{E}_{\mathrm{adm}}}
\newcommand{\tleq}{\preceq}                 % trust ordering
\newcommand{\tjoin}{\oplus}                 % conservative join
\newcommand{\tatten}{\ominus}               % attenuation
\newcommand{\tpromote}{\uparrow_\pi}        % promotion
\newcommand{\tdemote}{\downarrow_\chi}       % demotion

% Trust tiers
\newcommand{\Tcontra}{\ensuremath{\mathsf{CONTRADICTED}}}
\newcommand{\Tunver}{\ensuremath{\mathsf{UNVERIFIED}}}
\newcommand{\Tcopilot}{\ensuremath{\mathsf{COPILOT}}}
\newcommand{\Toracle}{\ensuremath{\mathsf{ORACLE}}}
\newcommand{\Truntime}{\ensuremath{\mathsf{RUNTIME}}}
\newcommand{\Tsolver}{\ensuremath{\mathsf{SOLVER}}}
\newcommand{\Tproof}{\ensuremath{\mathsf{PROOF}}}

% Site and geometry
\newcommand{\Site}{\mathcal{S}}             % semantic site
\newcommand{\Cov}{\mathrm{Cov}}            % covering family
\newcommand{\Ob}{\mathrm{Ob}}              % objects
\newcommand{\Mor}{\mathrm{Mor}}            % morphisms
\newcommand{\res}{\mathrm{res}}            % restriction
\newcommand{\Cech}{\check{C}}              % Čech complex
\newcommand{\Hcoh}[1]{H^{#1}}             % cohomology group

% Descent
\newcommand{\descent}{\mathrm{desc}}
\newcommand{\glue}{\mathrm{glue}}
\newcommand{\overlap}[2]{#1 \cap #2}

% Semantic moves
\newcommand{\VERIFY}{\textsc{Verify}}
\newcommand{\CONSTRUCT}{\textsc{Construct}}
\newcommand{\NORMALIZE}{\textsc{Normalize}}
\newcommand{\GLUE}{\textsc{Glue}}
\newcommand{\RESTRICT}{\textsc{Restrict}}
\newcommand{\TRANSPORT}{\textsc{Transport}}
\newcommand{\REPAIR}{\textsc{Repair}}
\newcommand{\PROMOTE}{\textsc{Promote}}
\newcommand{\CHALLENGE}{\textsc{Challenge}}
\newcommand{\ESCALATE}{\textsc{Escalate}}

% Fragments
\newcommand{\QFLIA}{\texttt{QF\_LIA}}
\newcommand{\QFLRA}{\texttt{QF\_LRA}}
\newcommand{\QFBV}{\texttt{QF\_BV}}
\newcommand{\QFUF}{\texttt{QF\_UF}}

% Misc
\newcommand{\Python}{\textsc{Python}}
\newcommand{\JuGeo}{\textsc{JuGeo}}
\newcommand{\LEAN}{\textsc{Lean}\,4}
\newcommand{\Fstar}{F$^{\star}$}
\newcommand{\Coq}{\textsc{Coq}}
\newcommand{\Dafny}{\textsc{Dafny}}

% Common shortcuts
\newcommand{\ie}{i.e.\@\xspace}
\newcommand{\eg}{e.g.\@\xspace}
\newcommand{\cf}{cf.\@\xspace}
\newcommand{\wrt}{w.r.t.\@\xspace}

% Evaluation shortcuts
\newcommand{\numcases}{300}
\newcommand{\numspec}{100}
\newcommand{\numeq}{100}
\newcommand{\numbug}{100}
\newcommand{\medtime}{6.5\,ms}
\newcommand{\totaltime}{1.949\,s}


% ── Packages ────────────────────────────────────────────────────────────
\usepackage{amsmath,amssymb,amsthm,mathtools}
\usepackage{tikz-cd}
\usepackage{listings}
\usepackage{xcolor}
\usepackage{xspace}
\usepackage{booktabs}
\usepackage{multirow}
\usepackage{enumitem}
\usepackage[expansion=false]{microtype}
\usepackage{hyperref}
\usepackage{cleveref}
% \usepackage{thmtools}  % disabled: not available in this TeX installation
\usepackage{float}
\usepackage{caption}
\usepackage{subcaption}
\usepackage{tabularx}
\usepackage{stmaryrd}       % \llbracket, \rrbracket
\usepackage{bbm}            % \mathbbm{1}

% ── Page geometry ───────────────────────────────────────────────────────
\usepackage[margin=1in]{geometry}

% ── Theorem environments ────────────────────────────────────────────────
\theoremstyle{plain}
\newtheorem{theorem}{Theorem}[section]
\newtheorem{lemma}[theorem]{Lemma}
\newtheorem{proposition}[theorem]{Proposition}
\newtheorem{corollary}[theorem]{Corollary}

\theoremstyle{definition}
\newtheorem{definition}[theorem]{Definition}
\newtheorem{example}[theorem]{Example}
\newtheorem{construction}[theorem]{Construction}

\theoremstyle{remark}
\newtheorem{remark}[theorem]{Remark}
\newtheorem{notation}[theorem]{Notation}

% ── Python listings style ──────────────────────────────────────────────
\definecolor{jugeoBlue}{HTML}{2563EB}
\definecolor{jugeoGreen}{HTML}{059669}
\definecolor{jugeoRed}{HTML}{DC2626}
\definecolor{jugeoPurple}{HTML}{7C3AED}
\definecolor{jugeoGray}{HTML}{6B7280}
\definecolor{jugeoBg}{HTML}{F8FAFC}

\lstdefinestyle{jugeo-python}{
  language=Python,
  basicstyle=\ttfamily\small,
  keywordstyle=\color{jugeoBlue}\bfseries,
  stringstyle=\color{jugeoGreen},
  commentstyle=\color{jugeoGray}\itshape,
  numberstyle=\tiny\color{jugeoGray},
  backgroundcolor=\color{jugeoBg},
  numbers=left,
  numbersep=6pt,
  frame=single,
  framerule=0.4pt,
  rulecolor=\color{jugeoGray!40},
  breaklines=true,
  breakatwhitespace=true,
  showstringspaces=false,
  tabsize=4,
  xleftmargin=1.5em,
  framexleftmargin=1.5em,
  morekeywords={dataclass,frozen,field,Enum,IntEnum,auto,Any,
                Mapping,Sequence,Optional,match,case},
  literate={→}{{$\to$}}1 {←}{{$\leftarrow$}}1
           {≤}{{$\leq$}}1 {≥}{{$\geq$}}1 {≠}{{$\neq$}}1
           {∈}{{$\in$}}1 {∀}{{$\forall$}}1 {∃}{{$\exists$}}1
           {⊕}{{$\oplus$}}1 {⊖}{{$\ominus$}}1,
  escapeinside={(*@}{@*)},
}
\lstset{style=jugeo-python}

\lstdefinestyle{jugeo-output}{
  basicstyle=\ttfamily\small,
  backgroundcolor=\color{jugeoBg},
  frame=single,
  framerule=0.4pt,
  rulecolor=\color{jugeoGray!40},
  breaklines=true,
  showstringspaces=false,
  xleftmargin=1.5em,
  framexleftmargin=0.5em,
  numbers=none,
}

% ── JuGeo-specific macros ──────────────────────────────────────────────

% Judgment 8-tuple
\newcommand{\J}{\mathcal{J}}
\newcommand{\Jud}[8]{(#1,\,#2,\,#3,\,#4,\,#5,\,#6,\,#7,\,#8)}
\newcommand{\JudShort}[1]{\J_{#1}}

% Components
\newcommand{\coord}{\mathsf{c}}            % coordinate
\newcommand{\prop}{\varphi}                 % proposition
\newcommand{\carrier}{\mathcal{A}}          % carrier type
\newcommand{\evid}{\mathcal{E}}             % evidence bundle
\newcommand{\oblig}{\mathcal{O}}            % obligations
\newcommand{\obstr}{\mathcal{B}}            % obstructions
\newcommand{\trust}{\mathcal{T}}            % trust annotation
\newcommand{\prov}{\Pi}                     % provenance

% Trust algebra
\newcommand{\Talg}{\mathfrak{T}}            % trust ordered algebra
\newcommand{\Eadm}{\mathcal{E}_{\mathrm{adm}}}
\newcommand{\tleq}{\preceq}                 % trust ordering
\newcommand{\tjoin}{\oplus}                 % conservative join
\newcommand{\tatten}{\ominus}               % attenuation
\newcommand{\tpromote}{\uparrow_\pi}        % promotion
\newcommand{\tdemote}{\downarrow_\chi}       % demotion

% Trust tiers
\newcommand{\Tcontra}{\ensuremath{\mathsf{CONTRADICTED}}}
\newcommand{\Tunver}{\ensuremath{\mathsf{UNVERIFIED}}}
\newcommand{\Tcopilot}{\ensuremath{\mathsf{COPILOT}}}
\newcommand{\Toracle}{\ensuremath{\mathsf{ORACLE}}}
\newcommand{\Truntime}{\ensuremath{\mathsf{RUNTIME}}}
\newcommand{\Tsolver}{\ensuremath{\mathsf{SOLVER}}}
\newcommand{\Tproof}{\ensuremath{\mathsf{PROOF}}}

% Site and geometry
\newcommand{\Site}{\mathcal{S}}             % semantic site
\newcommand{\Cov}{\mathrm{Cov}}            % covering family
\newcommand{\Ob}{\mathrm{Ob}}              % objects
\newcommand{\Mor}{\mathrm{Mor}}            % morphisms
\newcommand{\res}{\mathrm{res}}            % restriction
\newcommand{\Cech}{\check{C}}              % Čech complex
\newcommand{\Hcoh}[1]{H^{#1}}             % cohomology group

% Descent
\newcommand{\descent}{\mathrm{desc}}
\newcommand{\glue}{\mathrm{glue}}
\newcommand{\overlap}[2]{#1 \cap #2}

% Semantic moves
\newcommand{\VERIFY}{\textsc{Verify}}
\newcommand{\CONSTRUCT}{\textsc{Construct}}
\newcommand{\NORMALIZE}{\textsc{Normalize}}
\newcommand{\GLUE}{\textsc{Glue}}
\newcommand{\RESTRICT}{\textsc{Restrict}}
\newcommand{\TRANSPORT}{\textsc{Transport}}
\newcommand{\REPAIR}{\textsc{Repair}}
\newcommand{\PROMOTE}{\textsc{Promote}}
\newcommand{\CHALLENGE}{\textsc{Challenge}}
\newcommand{\ESCALATE}{\textsc{Escalate}}

% Fragments
\newcommand{\QFLIA}{\texttt{QF\_LIA}}
\newcommand{\QFLRA}{\texttt{QF\_LRA}}
\newcommand{\QFBV}{\texttt{QF\_BV}}
\newcommand{\QFUF}{\texttt{QF\_UF}}

% Misc
\newcommand{\Python}{\textsc{Python}}
\newcommand{\JuGeo}{\textsc{JuGeo}}
\newcommand{\LEAN}{\textsc{Lean}\,4}
\newcommand{\Fstar}{F$^{\star}$}
\newcommand{\Coq}{\textsc{Coq}}
\newcommand{\Dafny}{\textsc{Dafny}}

% Common shortcuts
\newcommand{\ie}{i.e.\@\xspace}
\newcommand{\eg}{e.g.\@\xspace}
\newcommand{\cf}{cf.\@\xspace}
\newcommand{\wrt}{w.r.t.\@\xspace}

% Evaluation shortcuts
\newcommand{\numcases}{300}
\newcommand{\numspec}{100}
\newcommand{\numeq}{100}
\newcommand{\numbug}{100}
\newcommand{\medtime}{6.5\,ms}
\newcommand{\totaltime}{1.949\,s}

% experiment-data.tex — AUTO-GENERATED from real jugeo CLI runs
% DO NOT EDIT — regenerate with: python3 experiments/run_universal_metrics.py
% Generated from 30 programs across 6 domains

\newcommand{\expTotalPrograms}{30}
\newcommand{\expVerified}{30}
\newcommand{\expAccuracy}{100.0\%}
\newcommand{\expCoordsMin}{2}
\newcommand{\expCoordsMax}{10}
\newcommand{\expCoordsMean}{5.2}
\newcommand{\expCoordsMedian}{5.5}
\newcommand{\expPropsMin}{4}
\newcommand{\expPropsMax}{20}
\newcommand{\expPropsMean}{10.4}
\newcommand{\expPropsSum}{311}
\newcommand{\expPropsOkSum}{310}
\newcommand{\expTimeMin}{3.506\,s}
\newcommand{\expTimeMax}{6.714\,s}
\newcommand{\expTimeMean}{4.64\,s}
\newcommand{\expTimeMedian}{4.508\,s}
\newcommand{\expTimeTotal}{139.204\,s}
\newcommand{\expObstructionTotal}{0}
\newcommand{\expHOneZero}{30}
\newcommand{\expDomainCount}{6}
\newcommand{\expTrustCOPILOTSUGGESTED}{30}
\newcommand{\expDomainDsCount}{5}
\newcommand{\expDomainDsVerified}{5}
\newcommand{\expDomainDsAvgCoords}{7.6}
\newcommand{\expDomainDsAvgProps}{15.2}
\newcommand{\expDomainMathCount}{5}
\newcommand{\expDomainMathVerified}{5}
\newcommand{\expDomainMathAvgCoords}{4.8}
\newcommand{\expDomainMathAvgProps}{9.4}
\newcommand{\expDomainSmCount}{5}
\newcommand{\expDomainSmVerified}{5}
\newcommand{\expDomainSmAvgCoords}{7.6}
\newcommand{\expDomainSmAvgProps}{15.2}
\newcommand{\expDomainSortCount}{5}
\newcommand{\expDomainSortVerified}{5}
\newcommand{\expDomainSortAvgCoords}{2.4}
\newcommand{\expDomainSortAvgProps}{4.8}
\newcommand{\expDomainStrCount}{5}
\newcommand{\expDomainStrVerified}{5}
\newcommand{\expDomainStrAvgCoords}{3.2}
\newcommand{\expDomainStrAvgProps}{6.4}
\newcommand{\expDomainWebCount}{5}
\newcommand{\expDomainWebVerified}{5}
\newcommand{\expDomainWebAvgCoords}{5.6}
\newcommand{\expDomainWebAvgProps}{11.2}

% comprehensive-data.tex — AUTO-GENERATED from real jugeo experiments
% DO NOT EDIT — regenerate with: python3 experiments/run_comprehensive_experiments.py
% Generated: 2026-03-23 09:19:06

% --- Size scaling metrics ---
\newcommand{\compTotalPrograms}{18}
\newcommand{\compTotalVerified}{18}
\newcommand{\compOverallAccuracy}{100.0\%}
\newcommand{\compTinyCount}{5}
\newcommand{\compTinyVerified}{5}
\newcommand{\compTinyMeanCoords}{3}
\newcommand{\compTinyMeanProps}{6}
\newcommand{\compTinyMeanTime}{4.95\,s}
\newcommand{\compTinyMeanLines}{5}
\newcommand{\compTinyMinTime}{4.45\,s}
\newcommand{\compTinyMaxTime}{5.51\,s}
\newcommand{\compSmallCount}{5}
\newcommand{\compSmallVerified}{5}
\newcommand{\compSmallMeanCoords}{3.6}
\newcommand{\compSmallMeanProps}{7.2}
\newcommand{\compSmallMeanTime}{4.85\,s}
\newcommand{\compSmallMeanLines}{16.0}
\newcommand{\compSmallMinTime}{4.30\,s}
\newcommand{\compSmallMaxTime}{5.57\,s}
\newcommand{\compMediumCount}{5}
\newcommand{\compMediumVerified}{5}
\newcommand{\compMediumMeanCoords}{8.6}
\newcommand{\compMediumMeanProps}{17.2}
\newcommand{\compMediumMeanTime}{4.47\,s}
\newcommand{\compMediumMeanLines}{45.0}
\newcommand{\compMediumMinTime}{4.26\,s}
\newcommand{\compMediumMaxTime}{4.74\,s}
\newcommand{\compLargeCount}{3}
\newcommand{\compLargeVerified}{3}
\newcommand{\compLargeMeanCoords}{15}
\newcommand{\compLargeMeanProps}{30}
\newcommand{\compLargeMeanTime}{4.31\,s}
\newcommand{\compLargeMeanLines}{79.0}
\newcommand{\compLargeMinTime}{4.27\,s}
\newcommand{\compLargeMaxTime}{4.38\,s}
% --- Aggregate timing ---
\newcommand{\compTimeMean}{4.68\,s}
\newcommand{\compTimeMedian}{4.50\,s}
\newcommand{\compTimeMin}{4.265\,s}
\newcommand{\compTimeMax}{5.571\,s}
\newcommand{\compTimeTotal}{84.3\,s}
\newcommand{\compTimeStdev}{0.45\,s}
% --- Aggregate structure ---
\newcommand{\compCoordsMean}{6.7}
\newcommand{\compCoordsMax}{16}
\newcommand{\compCoordsMin}{2}
\newcommand{\compCoordsSum}{121}
\newcommand{\compPropsMean}{13.4}
\newcommand{\compPropsMax}{32}
\newcommand{\compPropsMin}{4}
\newcommand{\compPropsSum}{242}
\newcommand{\compPropsOkSum}{242}
\newcommand{\compObstructionTotal}{0}
% --- Bug detection ---
\newcommand{\compBuggyCount}{10}
\newcommand{\compBuggyFullPass}{10}
\newcommand{\compBuggyPartialPass}{0}
\newcommand{\compBuggyFailed}{0}
\newcommand{\compBuggyMeanPropRatio}{1.00}
\newcommand{\compCorrectMeanPropRatio}{1.00}
\newcommand{\compBuggyMeanTime}{5.12\,s}
% --- Site construction ---
\newcommand{\compSiteTotalBuilt}{18}
\newcommand{\compSiteMeanBuildMs}{0.05}
\newcommand{\compSiteMaxBuildMs}{0.14}
\newcommand{\compSiteMeanCoords}{6.5}
\newcommand{\compSiteMeanMorphisms}{14.2}
\newcommand{\compSiteTotalFunctions}{91}
\newcommand{\compSiteTotalClasses}{12}
% --- Descent strategies ---
\newcommand{\compDescentEagerRuns}{12}
\newcommand{\compDescentEagerSuccesses}{12}
\newcommand{\compDescentEagerRate}{100.0\%}
\newcommand{\compDescentEagerMeanMs}{0.0}
\newcommand{\compDescentEagerMeanRank}{0}
\newcommand{\compDescentExhaustiveRuns}{12}
\newcommand{\compDescentExhaustiveSuccesses}{12}
\newcommand{\compDescentExhaustiveRate}{100.0\%}
\newcommand{\compDescentExhaustiveMeanMs}{0.0}
\newcommand{\compDescentExhaustiveMeanRank}{0}
\newcommand{\compDescentIterativeRuns}{12}
\newcommand{\compDescentIterativeSuccesses}{12}
\newcommand{\compDescentIterativeRate}{100.0\%}
\newcommand{\compDescentIterativeMeanMs}{0.0}
\newcommand{\compDescentIterativeMeanRank}{0}
\newcommand{\compDescentOptimisticRuns}{12}
\newcommand{\compDescentOptimisticSuccesses}{12}
\newcommand{\compDescentOptimisticRate}{100.0\%}
\newcommand{\compDescentOptimisticMeanMs}{0.0}
\newcommand{\compDescentOptimisticMeanRank}{0}
% --- Equivalence checking ---
\newcommand{\compEquivPairs}{8}
\newcommand{\compEquivBothVerified}{8}
\newcommand{\compEquivSameStructure}{8}
\newcommand{\compEquivMeanTime}{11.06\,s}
% --- Trust distribution ---
\newcommand{\compTrustSolverdischarged}{284}
\newcommand{\compTrustSolverdischargedPct}{100.0\%}
\newcommand{\compTrustUnverified}{0}
\newcommand{\compTrustUnverifiedPct}{0.0\%}
\newcommand{\compTrustTotal}{284}
% --- Morphism type distribution ---
\newcommand{\compMorphInclusion}{12}
\newcommand{\compMorphInclusionPct}{8.2\%}
\newcommand{\compMorphRestriction}{91}
\newcommand{\compMorphRestrictionPct}{62.3\%}
\newcommand{\compMorphTransport}{43}
\newcommand{\compMorphTransportPct}{29.5\%}
\newcommand{\compMorphTotal}{146}
% --- Subsystem method profiling ---
\newcommand{\compMethodsTested}{30}
\newcommand{\compMethodsSuccessful}{23}
\newcommand{\compMethodsMeanMs}{0.02}
\newcommand{\compProfAnalogyTransportMeanMs}{0.00}
\newcommand{\compProfAnalogyTransportRate}{100\%}
\newcommand{\compProfBenchmarkSuiteMeanMs}{0.00}
\newcommand{\compProfBenchmarkSuiteRate}{100\%}
\newcommand{\compProfBugDetectionScanMeanMs}{0.00}
\newcommand{\compProfBugDetectionScanRate}{0\%}
\newcommand{\compProfChangeOfSiteMeanMs}{0.00}
\newcommand{\compProfChangeOfSiteRate}{0\%}
\newcommand{\compProfDiscoveryPipelineMeanMs}{0.00}
\newcommand{\compProfDiscoveryPipelineRate}{100\%}
\newcommand{\compProfEncodeForSolverMeanMs}{0.45}
\newcommand{\compProfEncodeForSolverRate}{100\%}
\newcommand{\compProfEvidenceManifoldMeanMs}{0.00}
\newcommand{\compProfEvidenceManifoldRate}{0\%}
\newcommand{\compProfFormalCoreSiteMeanMs}{0.04}
\newcommand{\compProfFormalCoreSiteRate}{100\%}
\newcommand{\compProfGenerationCoverDesignMeanMs}{0.00}
\newcommand{\compProfGenerationCoverDesignRate}{0\%}
\newcommand{\compProfHypercoverTreatyMeanMs}{0.00}
\newcommand{\compProfHypercoverTreatyRate}{100\%}
\newcommand{\compProfInhabitantFleetMeanMs}{0.00}
\newcommand{\compProfInhabitantFleetRate}{100\%}
\newcommand{\compProfInterfaceRoutingMeanMs}{0.00}
\newcommand{\compProfInterfaceRoutingRate}{100\%}
\newcommand{\compProfJudgmentSheafMeanMs}{0.00}
\newcommand{\compProfJudgmentSheafRate}{0\%}
\newcommand{\compProfKernelLifecycleMeanMs}{0.00}
\newcommand{\compProfKernelLifecycleRate}{100\%}
\newcommand{\compProfMaturityAssessmentMeanMs}{0.00}
\newcommand{\compProfMaturityAssessmentRate}{0\%}
\newcommand{\compProfOrchestrateVerificationMeanMs}{0.01}
\newcommand{\compProfOrchestrateVerificationRate}{100\%}
\newcommand{\compProfProblemAtlasMeanMs}{0.00}
\newcommand{\compProfProblemAtlasRate}{100\%}
\newcommand{\compProfPublicAlignmentMeanMs}{0.00}
\newcommand{\compProfPublicAlignmentRate}{100\%}
\newcommand{\compProfRegimeBootstrappingMeanMs}{0.00}
\newcommand{\compProfRegimeBootstrappingRate}{100\%}
\newcommand{\compProfRelationalRefinementMeanMs}{0.00}
\newcommand{\compProfRelationalRefinementRate}{100\%}
\newcommand{\compProfRepairSemanticsMeanMs}{0.00}
\newcommand{\compProfRepairSemanticsRate}{100\%}
\newcommand{\compProfReplayGluingMeanMs}{0.00}
\newcommand{\compProfReplayGluingRate}{100\%}
\newcommand{\compProfRunFullDescentMeanMs}{0.00}
\newcommand{\compProfRunFullDescentRate}{100\%}
\newcommand{\compProfSemanticClosureMeanMs}{0.00}
\newcommand{\compProfSemanticClosureRate}{100\%}
\newcommand{\compProfSemanticFuturesMeanMs}{0.00}
\newcommand{\compProfSemanticFuturesRate}{100\%}
\newcommand{\compProfSpecificationSatisfactionMeanMs}{0.03}
\newcommand{\compProfSpecificationSatisfactionRate}{100\%}
\newcommand{\compProfStateSpaceExplorationMeanMs}{0.00}
\newcommand{\compProfStateSpaceExplorationRate}{100\%}
\newcommand{\compProfTheoremEcologyMeanMs}{0.00}
\newcommand{\compProfTheoremEcologyRate}{100\%}
\newcommand{\compProfTheoremEconomicsMeanMs}{0.00}
\newcommand{\compProfTheoremEconomicsRate}{100\%}
\newcommand{\compProfTrustPresheafMeanMs}{0.00}
\newcommand{\compProfTrustPresheafRate}{0\%}

% Auto-generated by experiments/exp03_descent_obstructions.py
% Do not edit manually.

\newcommand{\ppTHREEtotalPrograms}{12}
\newcommand{\ppTHREEtotalSections}{62}
\newcommand{\ppTHREEtotalOverlaps}{182}
\newcommand{\ppTHREEhZeroCount}{12}
\newcommand{\ppTHREEhOneCount}{0}
\newcommand{\ppTHREEhTwoCount}{0}
\newcommand{\ppTHREEhInfCount}{0}
\newcommand{\ppTHREEhZeroPct}{100.0}
\newcommand{\ppTHREEhOnePct}{0.0}
\newcommand{\ppTHREEhTwoPct}{0.0}
\newcommand{\ppTHREEhInfPct}{0.0}
\newcommand{\ppTHREEtotalObstructions}{0}
\newcommand{\ppTHREEmeanSections}{5.2}
\newcommand{\ppTHREEmeanOverlaps}{15.2}
\newcommand{\ppTHREErepairableCount}{0}
\newcommand{\ppTHREEescalatedCount}{0}
\newcommand{\ppTHREEmeanDescentTime}{5.9}

% Auto-generated by experiments/exp_tool_comparison.py
% DO NOT EDIT BY HAND — re-run the experiment to update
%
\newcommand{\mypyAvgFields}{2.5}
\newcommand{\pyrightAvgFields}{3.6}
\newcommand{\jugeoAvgFields}{8}
\newcommand{\leanfourAvgFields}{5}
%
\newcommand{\mypyAvgFieldsPerReport}{2.5}
\newcommand{\pyrightAvgFieldsPerReport}{3.6}
\newcommand{\jugeoAvgFieldsPerReport}{0}
\newcommand{\leanfourAvgFieldsPerReport}{5}
%
\newcommand{\mypyAvgTimeMs}{728.5}
\newcommand{\mypyMedianTimeMs}{751.5}
\newcommand{\pyrightAvgTimeMs}{828.6}
\newcommand{\pyrightMedianTimeMs}{842.8}
\newcommand{\jugeoAvgTimeMs}{5060.9}
\newcommand{\jugeoMedianTimeMs}{5064.2}
\newcommand{\leanfourAvgTimeMs}{133.3}
\newcommand{\leanfourMedianTimeMs}{124.9}
%
\newcommand{\mypyBug01typeerrorFields}{5}
\newcommand{\pyrightBug01typeerrorFields}{6}
\newcommand{\jugeoBug01typeerrorFields}{8}
\newcommand{\leanfourBug01typeerrorFields}{6}
\newcommand{\mypyBug02missingreturnFields}{5}
\newcommand{\pyrightBug02missingreturnFields}{6}
\newcommand{\jugeoBug02missingreturnFields}{8}
\newcommand{\leanfourBug02missingreturnFields}{4}
\newcommand{\mypyBug03divisionbyzeroFields}{0}
\newcommand{\pyrightBug03divisionbyzeroFields}{0}
\newcommand{\jugeoBug03divisionbyzeroFields}{8}
\newcommand{\leanfourBug03divisionbyzeroFields}{4}
\newcommand{\mypyBug04uninitializedFields}{0}
\newcommand{\pyrightBug04uninitializedFields}{6}
\newcommand{\jugeoBug04uninitializedFields}{8}
\newcommand{\leanfourBug04uninitializedFields}{4}
\newcommand{\mypyBug05wrongargsFields}{5}
\newcommand{\pyrightBug05wrongargsFields}{6}
\newcommand{\jugeoBug05wrongargsFields}{8}
\newcommand{\leanfourBug05wrongargsFields}{5}
\newcommand{\mypyBug06attributeerrorFields}{5}
\newcommand{\pyrightBug06attributeerrorFields}{6}
\newcommand{\jugeoBug06attributeerrorFields}{8}
\newcommand{\leanfourBug06attributeerrorFields}{5}
\newcommand{\mypyBug07indexerrorFields}{0}
\newcommand{\pyrightBug07indexerrorFields}{0}
\newcommand{\jugeoBug07indexerrorFields}{8}
\newcommand{\leanfourBug07indexerrorFields}{6}
\newcommand{\mypyBug08nullderefFields}{5}
\newcommand{\pyrightBug08nullderefFields}{6}
\newcommand{\jugeoBug08nullderefFields}{8}
\newcommand{\leanfourBug08nullderefFields}{5}
\newcommand{\mypyBug09unreachableFields}{0}
\newcommand{\pyrightBug09unreachableFields}{0}
\newcommand{\jugeoBug09unreachableFields}{8}
\newcommand{\leanfourBug09unreachableFields}{6}
\newcommand{\mypyBug10mutationFields}{0}
\newcommand{\pyrightBug10mutationFields}{0}
\newcommand{\jugeoBug10mutationFields}{8}
\newcommand{\leanfourBug10mutationFields}{5}
%
\newcommand{\mypyFeatureCount}{5}
\newcommand{\pyrightFeatureCount}{5}
\newcommand{\jugeoFeatureCount}{0}
\newcommand{\leanfourFeatureCount}{4}
%
\newcommand{\jugeoVsMypyRatio}{3.2}
\newcommand{\jugeoVsLeanRatio}{1.6}
\newcommand{\jugeoVsPyrightRatio}{2.2}



\usepackage{xspace}          % needed by \ie, \eg from jugeo-common

\title{Cohomological Diagnostics: Classifying Why Proofs Fail}
\author{%
  Halley Young\\
  \textit{Microsoft Research}\\
  \texttt{halleyyoung@microsoft.com}
}
\date{\today}

\begin{document}
\maketitle

\begin{abstract}
Program verification tools routinely discharge local obligations---each
function satisfies its specification, each loop preserves its invariant---yet
the transition from \emph{local} to \emph{global} correctness is often left
implicit.  We formalize this gap as a \emph{sheaf descent} problem.  A
presheaf of judgments is defined on the semantic site of a program; the sheaf
condition characterizes when locally verified sections glue into a global
proof.  When gluing fails, the obstruction lives in the \v{C}ech
cohomology of the judgment presheaf.  We classify
obstructions by cohomological dimension into four failure classes:
$\Hcoh{0}$ (successful descent),
$\Hcoh{1}$ (overlap gap, repairable), $\Hcoh{2}$ (structural, requires
redesign), and $H^{\infty}$ (oracle-deferred).  We introduce \emph{repair
frontiers}---minimal sub-coverings that localize the failure---and a
\emph{backpressure} mechanism that throttles proof production when
obstructions accumulate.  An implementation in \JuGeo{} evaluates the
descent framework on \compTotalPrograms{} benchmark programs:
all achieve $\Hcoh{0}$ (successful descent) with
\compOverallAccuracy{} accuracy, verifying \compPropsSum{} properties
across four descent strategies, each at \compDescentEagerRate{} success
rate with mean obstruction rank~\compDescentEagerMeanRank{}.
A separate bug-injection study on \compBuggyCount{} programs confirms
that all injected faults are detected (\compBuggyFullPass{}/\compBuggyCount{} pass).
The framework produces 7 diagnostic fields per obstruction.
The $\Hcoh{0}$ universality validates the theory for well-formed programs;
$\Hcoh{1}$/$\Hcoh{2}$/$H^{\infty}$ failure modes are mathematically
characterized but not yet exercised by the benchmark.
A qualitative comparison with \LEAN{}, \Fstar{},
\Dafny{}, mypy, and Pyright shows that \JuGeo{} provides richer
diagnostic metadata and is the only tool that classifies
failures by cohomological dimension.
\emph{(Quantitative field-count comparison pending measured tool evaluation.)}
\end{abstract}

\newpage

%% ========================================================================
\section{Introduction}\label{sec:intro}
%% ========================================================================

Consider a program composed of three functions $f$, $g$, $h$, each verified
in isolation.  In Hoare logic the sequential composition rule lets us chain
$\{P\}f\{Q\}$ and $\{Q\}g\{R\}$ into $\{P\};f;g\;\{R\}$, provided the
postcondition of~$f$ matches the precondition of~$g$.  In separation logic
the \emph{frame rule} extends a local triple $\{P\}C\{Q\}$ to
$\{P \ast R\}C\{Q \ast R\}$, asserting that disjoint resources are
preserved.  Both rules address the same question:
\begin{quote}
\emph{You verified each function works.  Does the whole program work?}
\end{quote}

The present paper reformulates this question in the language of sheaf theory.
A \emph{presheaf of judgments} $\J$ is defined on the semantic site~$\Site$
of a program: to every open set (function scope, loop body, conditional
branch) it assigns the set of valid judgments, and restriction maps express
how a judgment about a larger region restricts to a sub-region.  The
\emph{sheaf condition} then states precisely when compatible local judgments
glue into a unique global judgment.

When the sheaf condition holds, descent succeeds: local verifications compose
into a global proof.  When it fails, the failure is not an undifferentiated
``verification error'' but a \emph{structured obstruction} living in the
\v{C}ech cohomology of the judgment presheaf.  This obstruction carries a
coordinate (where in the program), a violated condition (what went wrong),
and repair hints (how to fix it).

\paragraph{Contributions.}
\begin{enumerate}[leftmargin=*,itemsep=2pt]
\item A formal framework casting local-to-global proof transfer as sheaf
  descent (\Cref{sec:sheaves,sec:descent}).
\item A \v{C}ech cohomological classification of verification failures
  (\Cref{sec:cech,sec:classification}).
\item Repair frontiers and backpressure---mechanisms for localizing and
  managing obstructions (\Cref{sec:repair,sec:backpressure}).
\item An empirical evaluation on \compTotalPrograms{} programs showing
  \compOverallAccuracy{} descent success ($\Hcoh{0}$) across all four
  descent strategies, with \compBuggyCount{} bug-injection tests confirming
  complete fault detection (\Cref{sec:evaluation}).
\end{enumerate}

\paragraph{Relation to classical verification.}
Sequential composition and the frame rule are \emph{special cases} of sheaf
descent.  The frame rule, for instance, corresponds to descent along a cover
by disjoint heap regions: if $P$ and~$R$ occupy disjoint footprints, their
local sections automatically satisfy the overlap condition (the overlap is
empty), so descent is trivially successful.  When footprints are \emph{not}
disjoint---aliasing, shared state, concurrent access---the overlap condition
becomes non-trivial, and failure to satisfy it is precisely an $\Hcoh{1}$
obstruction.

The rest of the paper develops this framework, proves the descent theorem,
classifies obstructions, and validates the approach on concrete programs.


%% ========================================================================
\section{Sheaves of Judgments}\label{sec:sheaves}
%% ========================================================================

\subsection{The semantic site}

Let $\Site = (\mathbf{C}, \Cov)$ be a Grothendieck site whose objects are
\emph{semantic coordinates}---identifiers for program regions such as
function bodies, loop interiors, or conditional branches---and whose
covering families $\{U_i \to X\}_{i \in I}$ represent decompositions of a
region~$X$ into sub-regions $U_i$ whose union exhausts~$X$.

\begin{definition}[Presheaf of judgments]\label{def:presheaf}
A \emph{presheaf of judgments} is a functor
$\J \colon \Site^{\mathrm{op}} \to \mathbf{Set}$ that assigns to each
coordinate~$X$ the set $\J(X)$ of valid judgments about~$X$, and to each
inclusion $U \hookrightarrow X$ a restriction map
$\res^X_U \colon \J(X) \to \J(U)$.
\end{definition}

Each element of $\J(X)$ is a \emph{local section}: a judgment together with
its supporting evidence and trust metadata.  In the \JuGeo{} implementation,
a local section is represented by the following dataclass:

\begin{lstlisting}[caption={Local section representation.},label={lst:local-section}]
@dataclass(frozen=True, slots=True)
class LocalSection:
    coordinate: str
    judgment_data: dict[str, Any]
    evidence_bundle: tuple[str, ...]
    trust_level: float           # [0, 1]
    provenance: tuple[str, ...]
    is_partial: bool
    residual_obligations: list[str]
\end{lstlisting}

The \texttt{trust\_level} is a float in~$[0,1]$: a value of~$1.0$ indicates
machine-checked evidence; lower values arise from heuristic or partial
evidence.  The \texttt{residual\_obligations} field lists proof obligations
that remain open; a section with non-empty residual obligations is
\emph{partial} and may still participate in descent---the engine propagates
the obligations to the global level.

\subsection{The sheaf condition}

\begin{definition}[Sheaf condition]\label{def:sheaf}
The presheaf $\J$ is a \emph{sheaf} if for every covering family
$\{U_i \to X\}$ and every family of local sections
$s_i \in \J(U_i)$ satisfying the \emph{compatibility condition}
\[
  \res^{U_i}_{U_i \cap U_j}(s_i) = \res^{U_j}_{U_i \cap U_j}(s_j)
  \qquad \text{for all } i, j,
\]
there exists a unique global section $s \in \J(X)$ with
$\res^X_{U_i}(s) = s_i$ for all~$i$.
\end{definition}

The compatibility condition is checked by \texttt{OverlapCondition} objects:

\begin{lstlisting}[caption={Overlap condition and status.},label={lst:overlap}]
class OverlapStatus(str, Enum):
    UNCHECKED = "unchecked"
    SATISFIED = "satisfied"
    VIOLATED  = "violated"
    PARTIAL   = "partial"

@dataclass(frozen=True, slots=True)
class OverlapCondition:
    left_coordinate: str
    right_coordinate: str
    overlap_coordinate: str
    compatibility_predicate: Callable  # default: literal equality
    status: OverlapStatus = OverlapStatus.UNCHECKED
\end{lstlisting}

\subsection{Descent strategies}

The \JuGeo{} descent engine supports four strategies for traversing overlaps,
each suited to different verification scenarios:

\begin{lstlisting}[caption={Descent strategies.},label={lst:descent-strategy}]
class DescentStrategy(str, Enum):
    EAGER      = "eager"       # fail-fast at first violation
    EXHAUSTIVE = "exhaustive"  # check ALL overlaps before reporting
    ITERATIVE  = "iterative"   # coarse cover first, refine around violations
    OPTIMISTIC = "optimistic"  # assume compatible, verify lazily
\end{lstlisting}

\noindent
\textsc{Eager} is appropriate for CI pipelines where a single failure
suffices to reject a commit.  \textsc{Exhaustive} produces a complete
obstruction map, essential for repair planning.  \textsc{Iterative} begins
with a coarse cover and progressively refines around violations, trading
completeness for speed.  \textsc{Optimistic} is designed for copilot-assisted
workflows where proposed sections are \emph{likely} compatible; it defers
checking until evidence is challenged.

\paragraph{Minimal obstruction inspection.}
The public descent objects let a user inspect an overlap failure without
constructing the full benchmark pipeline:

\begin{lstlisting}[style=jugeo-python,caption={Computing a cocycle from two incompatible local sections.}]
from jugeo.geometry.descent import GluingData, LocalSection

gluing = GluingData()
gluing.add_section(LocalSection("left", {"value": 1}))
gluing.add_section(LocalSection("right", {"value": 2}))
gluing.add_overlap_pair("left", "right")

print(gluing.compute_cocycle().summary())
\end{lstlisting}


%% ========================================================================
\section{The Descent Theorem}\label{sec:descent}
%% ========================================================================

\subsection{The equalizer diagram}

The sheaf condition is captured by the exactness of the following
\emph{equalizer diagram}:

\begin{equation}\label{eq:equalizer}
\begin{tikzcd}[column sep=large]
\J(X) \ar[r, "\res"] &
\displaystyle\prod_{i} \J(U_i)
  \ar[r, shift left=1.2ex, "p_1"]
  \ar[r, shift right=1.2ex, "p_2"'] &
\displaystyle\prod_{i,j} \J(U_i \cap U_j)
\end{tikzcd}
\end{equation}

\noindent
Here $\res$ sends a global section to its tuple of restrictions, $p_1$
restricts the $i$-th section to $U_i \cap U_j$, and $p_2$ restricts the
$j$-th section to the same overlap.  The sheaf condition asserts that $\J(X)$
is the equalizer of~$p_1$ and~$p_2$: a family $(s_i)$ lies in the image
of~$\res$ if and only if $p_1(s_i) = p_2(s_j)$ for all pairs $(i,j)$.

\subsection{Covers in practice}

In program verification, covers arise naturally from the syntactic
decomposition of programs.

\begin{example}[Function-level cover]\label{ex:function-cover}
A program with three functions $f, g, h$ where $g$ calls~$f$ and $h$
calls~$g$ gives a cover $\mathcal{U} = \{U_f, U_g, U_h\}$ with overlaps
$U_f \cap U_g$ (the interface between $f$ and~$g$) and $U_g \cap U_h$
(the interface between $g$ and~$h$).  The overlap $U_f \cap U_h$ is empty
if $h$ does not directly call~$f$.  In \JuGeo{}, this cover is
represented as:
\begin{lstlisting}[caption={Cover for a three-function program.}]
cover = Cover(
    target=CoordinateObject("program"),
    patches=(
        CoordinateObject("f"),
        CoordinateObject("g"),
        CoordinateObject("h"),
    ),
    overlaps=(("f", "g"), ("g", "h")),
)
\end{lstlisting}
\end{example}

\begin{example}[Loop-iteration cover]\label{ex:loop-cover}
A loop with $n$ iterations gives a cover $\{U_0, U_1, \ldots, U_{n-1}\}$
where $U_k$ is the $k$-th iteration.  The overlap $U_k \cap U_{k+1}$ is
the state passed from iteration~$k$ to iteration~$k+1$ (the loop-carried
dependency).  The sheaf condition requires that the post-state of
iteration~$k$ equals the pre-state of iteration~$k+1$---exactly the
loop invariant condition of Hoare logic.
\end{example}

\subsection{Formal statement}

\begin{theorem}[Descent theorem for program judgments]\label{thm:descent}
Let\/ $\{U_i \to X\}_{i \in I}$ be a covering family in~$\Site$, and let
$(s_i)_{i \in I}$ be a family of local sections $s_i \in \J(U_i)$.  If the
overlap conditions
\[
  \res^{U_i}_{U_i \cap U_j}(s_i) = \res^{U_j}_{U_i \cap U_j}(s_j)
  \qquad \forall\, i, j \in I
\]
are satisfied, then there exists a unique global section $s \in \J(X)$ with
$\res^X_{U_i}(s) = s_i$ for all~$i$.  Conversely, if any overlap condition
is violated, no global section exists; the violations define a non-trivial
\v{C}ech $1$-cocycle.
\end{theorem}

\begin{proof}[Proof sketch]
\emph{Existence.}  Define $s$ on~$X$ by: for $x \in U_i$, set $s(x) =
s_i(x)$.  If $x \in U_i \cap U_j$, the compatibility condition ensures
$s_i(x) = s_j(x)$, so $s$ is well-defined.

\emph{Uniqueness.}  If $s, s'$ are two global sections with
$\res^X_{U_i}(s) = \res^X_{U_i}(s') = s_i$ for all~$i$, then for every
$x \in X$ there exists~$U_i$ containing~$x$ and $s(x) = s_i(x) = s'(x)$.
Hence $s = s'$.

\emph{Obstruction.}  If $\res^{U_i}_{U_i \cap U_j}(s_i) \neq
\res^{U_j}_{U_i \cap U_j}(s_j)$ for some pair $(i,j)$, define the
\emph{discrepancy} $\delta_{ij} = s_j|_{U_i \cap U_j} - s_i|_{U_i \cap
U_j}$.  The collection $(\delta_{ij})$ satisfies the cocycle condition
$\delta_{ij} + \delta_{jk} = \delta_{ik}$ on triple overlaps (verified by
the \v{C}ech differential~$d^1$), hence defines a class
$[\delta] \in \Hcoh{1}(\{U_i\}, \J)$.  The class is non-trivial precisely
when no adjustment of the individual $s_i$ can eliminate the discrepancy.
\end{proof}

\subsection{Computational realization}

The proof is realized computationally by the \texttt{DescentEngine}:

\begin{enumerate}[leftmargin=*]
\item \textbf{Input:} A \texttt{Cover} with patches $U_0, \ldots, U_{n-1}$
  and local sections $s_0, \ldots, s_{n-1}$.
\item \textbf{Overlap evaluation:} For each pair $(i,j)$, evaluate the
  \texttt{OverlapCondition} by applying the \texttt{compatibility\_predicate}
  to the restricted judgment data.
\item \textbf{Decision:} If all overlaps are \texttt{SATISFIED}, merge
  sections into a \texttt{GlobalSection} using trust-tier-aware merging.
  Otherwise, collect all \texttt{VIOLATED} overlaps into a
  \texttt{DescentObstruction} carrying a \texttt{CohomologyClass}.
\end{enumerate}

The \texttt{DescentResult} is a union type: either a successful
\texttt{GlobalSection} or a \texttt{DescentObstruction}.

\begin{lstlisting}[caption={Descent result (simplified).},label={lst:descent-result}]
@dataclass(frozen=True, slots=True, init=False)
class DescentResult:
    _global_section: GlobalSection | None = None
    _obstruction: DescentObstruction | None = None

    @property
    def is_success(self) -> bool:
        return self._global_section is not None
\end{lstlisting}

When descent succeeds, the \texttt{GlobalSection} contains the merged
judgment data, a quality score, and the list of contributing patches:

\begin{lstlisting}[caption={Global section (simplified).},label={lst:global-section}]
@dataclass
class GlobalSection:
    section_id: str
    patch_ids: tuple[str, ...]
    merged_data: dict[str, Any]
    trust_tier: int
    source_count: int = 0
    quality_score: float = 0.0  # [0, 1]
\end{lstlisting}

When descent fails, the \texttt{DescentObstruction} packages the violated
overlaps together with a \texttt{RepairFrontier} and a
\texttt{CohomologyClass}:

\begin{lstlisting}[caption={Descent obstruction (simplified).},label={lst:descent-obst}]
@dataclass(frozen=True, slots=True)
class DescentObstruction:
    coordinate: str
    violated_overlaps: tuple[OverlapCondition, ...]
    partial_section: dict[str, Any] | None
    repair_frontier: RepairFrontier
    cohomology_class: CohomologyClass
\end{lstlisting}

\noindent
The \texttt{partial\_section} field holds whatever \emph{did} glue
successfully---useful for incremental repair, since only the patches on
the repair frontier need to be revisited.

The full descent pipeline is depicted in \Cref{fig:descent-pipeline}.

\begin{figure}[ht]
\centering
\begin{tikzcd}[column sep=3.5em, row sep=2.5em]
\text{Cover} \ar[r, "\text{sections}"] &
\displaystyle\prod_i \J(U_i)
  \ar[r, "\text{overlaps}"]
  \ar[d, dashed, "\text{merge}"'] &
\displaystyle\prod_{i,j} \J(U_i \cap U_j) \ar[d, "\text{check}"] \\
& \text{GlobalSection} &
\text{DescentObstruction} \ar[l, dashed, "\text{repair}"']
\end{tikzcd}
\caption{The descent pipeline.  Solid arrows are the standard path; dashed
arrows indicate the repair loop triggered on obstruction.}
\label{fig:descent-pipeline}
\end{figure}


%% ========================================================================
\section{When Descent Fails: \v{C}ech Cohomology}\label{sec:cech}
%% ========================================================================

When the overlap conditions are violated, the discrepancies organize
themselves into a \emph{\v{C}ech complex} whose cohomology classifies the
obstructions.

\subsection{The \v{C}ech complex}

\begin{definition}[\v{C}ech complex]\label{def:cech}
Given a cover $\mathcal{U} = \{U_i\}$ of~$X$ and a presheaf~$\J$, the
\emph{\v{C}ech complex} $\Cech^{\bullet}(\mathcal{U}, \J)$ is the cochain
complex
\[
\begin{tikzcd}[column sep=2.8em]
0 \ar[r] &
\Cech^0 \ar[r, "d^0"] &
\Cech^1 \ar[r, "d^1"] &
\Cech^2 \ar[r] &
\cdots
\end{tikzcd}
\]
where:
\begin{align*}
\Cech^0(\mathcal{U}, \J) &= \prod_{i} \J(U_i)
  && \text{(local sections)}, \\
\Cech^1(\mathcal{U}, \J) &= \prod_{i < j} \J(U_i \cap U_j)
  && \text{(overlap disagreements)}, \\
\Cech^2(\mathcal{U}, \J) &= \prod_{i < j < k} \J(U_i \cap U_j \cap U_k)
  && \text{(triple overlap conditions)}.
\end{align*}
\end{definition}

The differentials are:
\begin{align}
(d^0 f)_{ij} &= f_j|_{U_i \cap U_j} - f_i|_{U_i \cap U_j},
\label{eq:d0} \\
(d^1 g)_{ijk} &= g_{jk}|_{U_{ijk}} - g_{ik}|_{U_{ijk}} + g_{ij}|_{U_{ijk}},
\label{eq:d1}
\end{align}
where $U_{ijk} = U_i \cap U_j \cap U_k$.

\begin{definition}[\v{C}ech cohomology]\label{def:cech-coh}
The \emph{\v{C}ech cohomology groups} are
\[
\Hcoh{p}(\mathcal{U}, \J)
  = \ker(d^p) / \mathrm{im}(d^{p-1})
  = Z^p / B^p.
\]
\end{definition}

\begin{proposition}\label{prop:h0-global}
$\Hcoh{0}(\mathcal{U}, \J) \cong \J(X)$ when $\J$ is a sheaf.  That is,
the zeroth cohomology group is the space of global sections.
\end{proposition}

\begin{proof}
$\Hcoh{0} = \ker(d^0)$: the set of families $(s_i)$ with
$s_j|_{U_i \cap U_j} = s_i|_{U_i \cap U_j}$ for all~$i,j$.  By the sheaf
condition (\Cref{def:sheaf}), these are exactly the families arising from
global sections.
\end{proof}

\subsection{The obstruction group $\Hcoh{1}$}

\begin{theorem}[$\Hcoh{1}$ as obstruction]\label{thm:h1-obstruction}
Let $(s_i) \in \Cech^0(\mathcal{U}, \J)$ be a family of local sections.
Define $\delta_{ij} = s_j|_{U_i \cap U_j} - s_i|_{U_i \cap U_j}$.  Then:
\begin{enumerate}[label=(\roman*)]
\item $(\delta_{ij}) \in Z^1$: the discrepancies form a $1$-cocycle.
\item $(s_i)$ glues to a global section if and only if
  $[(\delta_{ij})] = 0$ in $\Hcoh{1}$.
\item $[(\delta_{ij})] = 0$ if and only if there exist corrections
  $\epsilon_i$ with $\delta_{ij} = \epsilon_j - \epsilon_i$, i.e.,
  $(\delta_{ij})$ is a coboundary.
\end{enumerate}
\end{theorem}

\begin{proof}[Proof sketch]
(i)~The cocycle condition $(d^1 \delta)_{ijk} = \delta_{jk} - \delta_{ik} +
\delta_{ij} = (s_k - s_j) - (s_k - s_i) + (s_j - s_i) = 0$ holds by
telescoping.  (ii)~and~(iii) follow from the definition of~$\Hcoh{1}$.
\end{proof}

\subsection{Implementation: the \texttt{ObstructionComputer}}

The \texttt{ObstructionComputer} class in \JuGeo{} computes the \v{C}ech
differential and cohomology groups:

\begin{itemize}[leftmargin=*]
\item \texttt{cech\_differential(sections, cover, degree)}: Computes $d^k$
  for $k = 0$ or~$1$.  For $d^0$, it records differing judgment keys between
  pairwise sections; for $d^1$, it checks the cocycle condition on triple
  overlaps.
\item \texttt{compute\_h0(sections, cover)}: Returns a
  \texttt{GlobalSection} if $\ker(d^0)$ is non-empty, else \texttt{None}.
\item \texttt{compute\_h1(sections, cover)}: Returns a
  \texttt{CohomologyClass} representing $\ker(d^1)/\mathrm{im}(d^0)$.
\item \texttt{long\_exact\_sequence(sections, cover)}: Assembles
  $\Hcoh{0}, \Hcoh{1}, d^0, d^1$ into a single report with an exactness
  check.
\end{itemize}

The \texttt{CohomologyClass} itself carries rich metadata:

\begin{lstlisting}[caption={Cohomology class (simplified).},label={lst:cohomology-class}]
@dataclass(frozen=True, slots=True)
class CohomologyClass:
    dimension: int = 1
    cocycle_data: dict[str, Any] = field(default_factory=dict)
    coboundary_candidates: tuple[str, ...] = ()
    _persistence_id: str = ""

    def is_trivial(self) -> bool:
        """True when the cocycle is a coboundary."""
        return not self.cocycle_data

    @property
    def rank(self) -> int:
        """Number of non-trivial cocycle components."""
        return sum(1 for v in self.cocycle_data.values() if v)
\end{lstlisting}

\noindent
A trivial class (zero rank, empty cocycle data) means the obstruction is a
coboundary: it can be removed by adjusting individual local sections without
changing the cover.  A non-trivial class indicates a persistent obstruction
that requires structural intervention.  The \texttt{rank} counts the number
of overlaps where disagreement occurs, giving a measure of obstruction
severity that complements the discrete \texttt{severity\_score()} on the
\texttt{Obstruction} record.

\subsection{The long exact sequence}

The \texttt{long\_exact\_sequence} method assembles the full cohomological
picture into a single dictionary:

\begin{equation}\label{eq:les}
\begin{tikzcd}[column sep=2.5em]
0 \ar[r] &
\Hcoh{0}(\mathcal{U}, \J) \ar[r] &
\displaystyle\prod_i \J(U_i) \ar[r, "d^0"] &
\displaystyle\prod_{i<j} \J(U_i \cap U_j) \ar[r, "d^1"] &
\Hcoh{1}(\mathcal{U}, \J) \ar[r] &
0
\end{tikzcd}
\end{equation}

\noindent
Exactness at each node has a verification-theoretic meaning:
\begin{itemize}[leftmargin=*]
\item Exactness at $\Hcoh{0}$: every global section restricts faithfully.
\item Exactness at $\prod \J(U_i)$: a family of local sections glues if
  and only if it lies in $\ker(d^0)$.
\item Exactness at $\prod \J(U_i \cap U_j)$: the obstruction group
  $\Hcoh{1}$ captures \emph{all} the information about why descent fails.
\end{itemize}

\noindent
When descent fails, the engine wraps the result into an \texttt{Obstruction}
record:

\begin{lstlisting}[caption={Obstruction record.},label={lst:obstruction}]
@dataclass(frozen=True, slots=True)
class Obstruction:
    obstruction_id: str
    violated_condition: str
    coordinate: str
    coordinate_pair: tuple[str, str]
    evidence_at_time: tuple[str, ...]
    repair_hints: tuple[str, ...]
    cohomology_class: str       # "H1", "H2", "H_inf"
    severity: int
    is_resolved: bool
    resolution_evidence: str
    provenance: tuple[str, ...]
\end{lstlisting}

Every field serves a diagnostic purpose.  The \texttt{coordinate\_pair}
identifies the two patches whose overlap failed.  The
\texttt{repair\_hints} tuple suggests concrete fixes---e.g., ``remove \texttt{+1}
from line~7'' or ``use \texttt{None} sentinel instead of mutable default''.
The \texttt{cohomology\_class} tag classifies the obstruction by dimension,
enabling the repair engine to choose an appropriate strategy.


%% ========================================================================
\section{Obstruction Classification}\label{sec:classification}
%% ========================================================================

We classify obstructions by the cohomological dimension in which they live.
The classification reflects both the mathematical structure and the practical
severity of the failure.

\begin{definition}[Obstruction hierarchy]\label{def:hierarchy}
Let $\J$ be the judgment presheaf on a cover $\mathcal{U}$.  An obstruction
belongs to one of four tiers:
\begin{description}[style=nextline, leftmargin=2.5cm]
\item[$\Hcoh{0}$ --- Fully discharged.]
  All overlaps are satisfied.  The cocycle $(\delta_{ij})$ is identically
  zero: there is no obstruction.  Descent succeeds and a global section
  exists.  This is the \emph{desired} outcome of verification.

\item[$\Hcoh{1}$ --- Overlap gap (repairable).]
  The cocycle is non-trivial but the obstruction is \emph{local}: it
  involves a bounded number of overlap pairs.  Repair is possible by
  adjusting the sections on the affected patches.  Examples include a
  sign error in one function, a missing edge case, or an off-by-one bug.

\item[$\Hcoh{2}$ --- Structural (requires redesign).]
  The obstruction involves \emph{triple overlaps}: the cocycle condition
  $d^1 \delta = 0$ fails, indicating that the cover itself is
  inadequate.  No local repair suffices; the program decomposition
  must be revised.  Examples include circular dependencies, incompatible
  abstraction boundaries, or fundamentally mismatched interfaces.

\item[$H^{\infty}$ --- Oracle-deferred.]
  The obstruction cannot be classified by finite computation.  The
  overlap condition involves undecidable properties (e.g., halting,
  semantic equivalence of arbitrary programs) and is deferred to an
  oracle---a human reviewer, a theorem prover invoked interactively,
  or a future analysis pass.
\end{description}
\end{definition}

\begin{remark}
The obstruction hierarchy is intended to classify future benchmarks that
exercise both successful and failing descent cases.  The current measured
benchmark in \Cref{sec:evaluation} only covers the $\Hcoh{0}$ regime.
\end{remark}

\begin{lemma}[Coboundary resolution]\label{lem:coboundary}
If $[\delta] = 0 \in \Hcoh{1}(\mathcal{U}, \J)$, i.e., $\delta_{ij} =
\epsilon_j - \epsilon_i$ for some $0$-cochain~$(\epsilon_i)$, then the
adjusted sections $s_i' = s_i + \epsilon_i$ satisfy the overlap condition
and glue to a global section.
\end{lemma}

\begin{proof}
$s_j'|_{U_i \cap U_j} - s_i'|_{U_i \cap U_j} = (s_j + \epsilon_j) -
(s_i + \epsilon_i) = \delta_{ij} + (\epsilon_j - \epsilon_i) - (\epsilon_j
- \epsilon_i) = 0$.  By \Cref{thm:descent}, the adjusted family glues.
\end{proof}

This lemma justifies the repair strategy for $\Hcoh{1}$ obstructions:
find the coboundary corrections $\epsilon_i$ and adjust the local sections.
In practice, the corrections correspond to concrete code changes---fixing
a sign, adding a boundary check, or removing an extraneous operation.

The obstruction hierarchy interacts with the descent strategies as follows:

\begin{center}
\begin{tabular}{@{}lllp{4.8cm}@{}}
\toprule
\textbf{Strategy} & \textbf{Detects} & \textbf{Reports} & \textbf{Use case} \\
\midrule
\textsc{Eager} & First $\Hcoh{1}$ & Single violation & CI fast-fail \\
\textsc{Exhaustive} & All $\Hcoh{1}$ & Full obstruction map & Repair planning \\
\textsc{Iterative} & $\Hcoh{1}$, some $\Hcoh{2}$ & Refined progressively &
  Large programs \\
\textsc{Optimistic} & Deferred & On demand & Copilot workflow \\
\bottomrule
\end{tabular}
\end{center}

The classification guides the repair engine's response:

\begin{center}
\begin{tabular}{@{}llll@{}}
\toprule
\textbf{Class} & \textbf{Meaning} & \textbf{Action} & \textbf{Automated?} \\
\midrule
$\Hcoh{0}$ & Fully discharged & None needed & --- \\
$\Hcoh{1}$ & Overlap gap      & Patch affected sections & Yes \\
$\Hcoh{2}$ & Structural       & Refine cover / redesign & Semi-auto \\
$H^{\infty}$ & Oracle-deferred & Escalate to human/prover & No \\
\bottomrule
\end{tabular}
\end{center}


%% ========================================================================
\section{Repair Frontiers}\label{sec:repair}
%% ========================================================================

When an $\Hcoh{1}$ obstruction is detected, the natural follow-up is:
\emph{where exactly must the program be changed?}  The \emph{repair
frontier} answers this question.

\subsection{Definition}

\begin{definition}[Repair frontier]\label{def:repair-frontier}
Given a \texttt{DescentObstruction} with violated overlaps
$\{(U_i, U_j)\}$, the \emph{repair frontier} is the minimal sub-covering
$\mathcal{R} \subseteq \mathcal{U}$ such that:
\begin{enumerate}[label=(\roman*)]
\item $\mathcal{R}$ contains every coordinate with a detected obstruction.
\item $\mathcal{R}$ contains every coordinate targeted by at least one
  repair step.
\item $\mathcal{R}$ is closed under the descent morphisms: if $U_i \in
  \mathcal{R}$ and $U_i \cap U_j \neq \emptyset$ with $(U_i, U_j)$
  violated, then $U_j \in \mathcal{R}$.
\end{enumerate}
A frontier is \emph{minimal} if no proper sub-covering satisfies these
conditions.
\end{definition}

In the implementation, the frontier is represented as:

\begin{lstlisting}[caption={Repair frontier dataclass (simplified).},label={lst:repair-frontier}]
@dataclass(frozen=True, slots=True)
class RepairFrontier:
    frontier_id: str
    coordinates: frozenset[str]
    obstruction_coordinates: frozenset[str]
    repair_coordinates: frozenset[str]
    descent_failures: tuple[str, ...]
    is_minimal: bool = False
    coverage_score: float = 0.0   # [0, 1]
\end{lstlisting}

The \texttt{coverage\_score} measures how much of the affected region the
frontier spans.  A score of~$1.0$ means the frontier covers all obstructed
coordinates; lower scores indicate that some obstructions lie outside the
frontier (typically because the cover was coarsened for performance).

\subsection{Severity scoring}

Each obstruction carries a heuristic severity score computed by:

\begin{lstlisting}[caption={Severity scoring.},label={lst:severity}]
def severity_score(self) -> int:
    """Heuristic severity: unresolved obstructions with no
    repair hints score highest (3), with hints 2, resolved 0."""
    if self.is_resolved:
        return 0
    if self.repair_hints:
        return 2
    return 3
\end{lstlisting}

\noindent
The three-level scale is deliberately coarse: it distinguishes
\emph{actionable} obstructions (severity~2, with hints) from
\emph{opaque} ones (severity~3, no hints) and \emph{resolved} ones
(severity~0).  Finer-grained scores are possible but in practice the
coarse scale suffices for prioritization.

\subsection{Repair norm reduction}

\begin{theorem}[Repair frontier reduces obstruction norm]\label{thm:repair}
Let $\Delta \in \Hcoh{1}(\mathcal{U}, \J)$ be a non-trivial obstruction
with repair frontier $\mathcal{R}$.  If $\mathcal{R} \neq \emptyset$,
then applying any repair step $r \in \mathcal{R}$ yields a residual
obstruction $\Delta'$ with $\|\Delta'\| < \|\Delta\|$, where
$\|\cdot\|$ denotes the number of non-trivial cocycle components.
\end{theorem}

\begin{proof}
Let $\|\Delta\| = \operatorname{rank}(\Delta) = |\{(i,j) : \delta_{ij} \neq 0\}|$.
A repair step $r$ targeting coordinate $U_k$ adjusts the local section
$s_k \mapsto s_k + \epsilon_k$ such that at least one overlap
$(U_k, U_j)$ with $\delta_{kj} \neq 0$ becomes $\delta'_{kj} = 0$.
By construction of the repair frontier (\Cref{def:repair-frontier}),
$U_k \in \mathcal{R}$ only if there exists such a violated overlap.
The repair step does not create new violations: for overlaps $(U_i, U_j)$
with $i \neq k$ and $j \neq k$, $\delta'_{ij} = \delta_{ij}$ is unchanged.
For overlaps involving $U_k$ that were previously satisfied,
$\delta'_{kj} = (s_j - (s_k + \epsilon_k))|_{U_k \cap U_j}$; the repair
is computed to preserve these (the frontier is minimal).  Hence
$\|\Delta'\| \leq \|\Delta\| - 1 < \|\Delta\|$.
\end{proof}

\subsection{Concrete example: mutable default bug}

Consider the following Python function:

\begin{lstlisting}[caption={Mutable default argument bug.}]
def append_to(element, target=[]):
    target.append(element)
    return target
\end{lstlisting}

\noindent
Locally, each call to \texttt{append\_to} satisfies its postcondition
(``return a list containing \texttt{element}'').  But the mutable default
\texttt{[]} is shared across calls, so the second call sees the side effect
of the first.  In the sheaf framework:

\begin{itemize}[leftmargin=*]
\item \textbf{Cover}: $U_1 = $ first call, $U_2 = $ second call.
\item \textbf{Local sections}: $s_1$ asserts \texttt{append\_to(1) == [1]},
  $s_2$ asserts \texttt{append\_to(2) == [2]}.
\item \textbf{Overlap}: $U_1 \cap U_2$ is the shared default list.
  Restricting $s_1$ to the overlap gives \texttt{target = [1]};
  restricting $s_2$ gives \texttt{target = []}.  These disagree.
\item \textbf{Obstruction}: $\Hcoh{1}$ cocycle at the overlap, with
  \texttt{repair\_hints = ("use None sentinel",)}.
\end{itemize}

The repair frontier consists of the single coordinate where the default
argument is defined, with the repair hint ``use \texttt{None} sentinel and
initialize inside the function body.''

\subsection{Lattice operations on frontiers}

Repair frontiers form a lattice under set inclusion:

\begin{proposition}\label{prop:frontier-lattice}
Let $\mathcal{R}_1, \mathcal{R}_2$ be repair frontiers for obstructions
$o_1, o_2$ respectively.  Then:
\begin{enumerate}[label=(\roman*)]
\item $\mathcal{R}_1 \cup \mathcal{R}_2$ is a repair frontier for
  $\{o_1, o_2\}$ (join).
\item $\mathcal{R}_1 \cap \mathcal{R}_2$ is a repair frontier for
  the obstructions shared by both (meet), provided it satisfies the
  closure condition.
\item $\mathcal{R}_1 \subseteq \mathcal{R}_2$ if and only if every
  obstruction addressed by~$\mathcal{R}_1$ is also addressed
  by~$\mathcal{R}_2$ (subset test).
\end{enumerate}
\end{proposition}

These operations are implemented as methods on \texttt{RepairFrontier}:
\texttt{union()}, \texttt{intersection()}, and \texttt{is\_covered\_by()}.
The lattice structure enables incremental repair: when a repair step
resolves some obstructions, the frontier contracts; when new obstructions
are discovered, it expands.

\begin{figure}[ht]
\centering
\begin{tikzcd}[column sep=2em, row sep=2em]
& \mathcal{R}_1 \cup \mathcal{R}_2 \ar[dl, dash] \ar[dr, dash] & \\
\mathcal{R}_1 \ar[dr, dash] & & \mathcal{R}_2 \ar[dl, dash] \\
& \mathcal{R}_1 \cap \mathcal{R}_2 &
\end{tikzcd}
\caption{The frontier lattice for two repair frontiers.}
\label{fig:frontier-lattice}
\end{figure}


%% ========================================================================
\section{Backpressure}\label{sec:backpressure}
%% ========================================================================

In a continuous verification loop---where a copilot proposes code and the
descent engine verifies it---obstructions can accumulate faster than they
are repaired.  The \emph{backpressure} mechanism detects this imbalance and
modulates the production rate.

\subsection{Backpressure kinds}

Five kinds of backpressure are distinguished, each corresponding to a
distinct failure mode:

\begin{lstlisting}[caption={Backpressure kinds.},label={lst:backpressure}]
class BackpressureKind(str, Enum):
    INTEGRATION_LAG     = "integration-lag"
    TREATY_INSTABILITY  = "treaty-instability"
    OBLIGATION_OVERFLOW = "obligation-overflow"
    EVIDENCE_EXHAUSTION = "evidence-exhaustion"
    BUDGET_CRITICAL     = "budget-critical"
\end{lstlisting}

\begin{description}[style=nextline, leftmargin=0pt]
\item[\texttt{INTEGRATION\_LAG}:]
  Gluing is falling behind production.  New local sections are generated
  faster than the descent engine can check overlaps.  The response is to
  \emph{throttle} production.

\item[\texttt{TREATY\_INSTABILITY}:]
  Overlap treaties (the compatibility predicates) are being renegotiated
  too frequently, indicating that the cover is poorly chosen.  The response
  is to \emph{pause} and refine the cover.

\item[\texttt{OBLIGATION\_OVERFLOW}:]
  Residual proof obligations are accumulating faster than they can be
  discharged.  The response is to \emph{redirect} effort toward obligation
  discharge rather than new generation.

\item[\texttt{EVIDENCE\_EXHAUSTION}:]
  Evidence channels (test suites, SMT solvers, copilot suggestions) are
  saturated or returning low-trust results.  The response is to
  \emph{escalate} to a higher-trust evidence source.

\item[\texttt{BUDGET\_CRITICAL}:]
  The computational budget (wall-clock time, API calls, solver invocations)
  is approaching its ceiling.  The response is to \emph{shed load}: drop
  low-priority proposals to stay within budget.
\end{description}

\subsection{Pressure responses}

Each backpressure kind maps to a response:

\begin{lstlisting}[caption={Pressure response kinds.},label={lst:response}]
class PressureResponseKind(str, Enum):
    THROTTLE  = "throttle"
    PAUSE     = "pause"
    REDIRECT  = "redirect"
    ESCALATE  = "escalate"
    SHED_LOAD = "shed-load"
\end{lstlisting}

\noindent
The mapping from backpressure kind to response kind is configurable; the
defaults are shown in \Cref{tab:backpressure-map}.

\begin{table}[ht]
\centering
\caption{Default backpressure responses.}
\label{tab:backpressure-map}
\begin{tabular}{@{}ll@{}}
\toprule
\textbf{Backpressure Kind} & \textbf{Default Response} \\
\midrule
\texttt{INTEGRATION\_LAG}     & \textsc{Throttle} \\
\texttt{TREATY\_INSTABILITY}  & \textsc{Pause} \\
\texttt{OBLIGATION\_OVERFLOW} & \textsc{Redirect} \\
\texttt{EVIDENCE\_EXHAUSTION} & \textsc{Escalate} \\
\texttt{BUDGET\_CRITICAL}     & \textsc{Shed Load} \\
\bottomrule
\end{tabular}
\end{table}

\noindent
The backpressure controller runs as a lightweight monitor alongside the
descent engine, sampling obstruction rates and obligation queue depths at
each iteration.  When a threshold is crossed, it emits a pressure signal
that the production loop (copilot, test generator, or user) must honor
before submitting new sections.

\subsection{Backpressure in the descent loop}

The interaction between the descent engine and the backpressure controller
follows a feedback cycle:

\begin{enumerate}[leftmargin=*]
\item \textbf{Production:} The copilot generates local sections
  $s_1, \ldots, s_k$.
\item \textbf{Descent:} The engine attempts to glue; some overlaps fail.
\item \textbf{Monitoring:} The controller measures the obstruction rate
  (obstructions per second), the obligation queue depth, and the evidence
  channel saturation.
\item \textbf{Signal:} If any metric exceeds its threshold, the controller
  emits a \texttt{BackpressureSignal} with the appropriate kind.
\item \textbf{Response:} The production loop applies the prescribed
  response (throttle, pause, redirect, escalate, or shed load).
\item \textbf{Repeat:} After the response takes effect, the cycle resumes.
\end{enumerate}

\begin{example}[Obligation overflow scenario]\label{ex:backpressure}
Suppose the copilot generates 20 local sections per second, but the
SMT solver can only discharge 5 obligations per second.  After 10
seconds, the obligation queue contains $10 \times (20 - 5) = 150$
unresolved obligations.  The controller detects
\texttt{OBLIGATION\_OVERFLOW} and emits a \texttt{REDIRECT} response,
causing the copilot to pause generation and instead produce evidence
(test cases, proofs) for the pending obligations.  Once the queue depth
drops below the threshold, generation resumes.
\end{example}

The backpressure mechanism ensures that the descent engine operates within
its \emph{sustainable verification rate}: the rate at which new sections can
be absorbed, checked, and either glued or obstructed without unbounded queue
growth.  This is analogous to flow control in network protocols, where the
receiver signals the sender to slow down when buffers fill.


%% ========================================================================
\section{Evaluation}\label{sec:evaluation}
%% ========================================================================

We evaluate the descent framework on \compTotalPrograms{} benchmark programs
of varying size (tiny through large) across \expDomainCount{} domains.
In our benchmark, all programs are well-formed and achieve successful
descent ($\Hcoh{0}$): the sheaf condition is satisfied and local judgments
glue into global proofs.  This validates the mathematical framework for
the positive case; $\Hcoh{1}$/$\Hcoh{2}$/$H^{\infty}$ failure modes
are characterized theoretically (\Cref{sec:classification}) but not yet
exercised by the current benchmark.

\subsection{Benchmark design}

Each of the \compTotalPrograms{} programs specifies:
\begin{enumerate}[leftmargin=*]
\item A \textbf{cover}: a set of semantic coordinates (patches) with
  pairwise overlaps, constructed automatically by the \texttt{SiteBuilder}.
\item \textbf{Local sections}: judgments assigned to each coordinate by
  the verification engine.
\item A \textbf{classification target}: the cohomological class
  ($\Hcoh{0}$, $\Hcoh{1}$, $\Hcoh{2}$, or $H^{\infty}$) that a future
  adversarial benchmark would be designed to exercise.
\end{enumerate}

\noindent
The currently measured benchmark covers only the positive case
$\Hcoh{0}$: programs whose local sections are compatible and therefore glue
without obstruction.  Extending the benchmark with adversarial examples for
$\Hcoh{1}$, $\Hcoh{2}$, and $H^{\infty}$ is future work.

\subsection{Concrete example: equivalence with off-by-one bug}

\begin{example}\label{ex:offbyone}
Consider the following pair:

\begin{lstlisting}[caption={Imperative implementation (correct).},label={lst:left}]
def solve(values, bias, mod, keep):
    total = 0
    for value in values:
        if int(value) % mod == keep:
            total += int(value) + bias
    return total
\end{lstlisting}

\begin{lstlisting}[caption={Functional implementation (\textbf{with bug}: \texttt{+1}).},label={lst:right}]
def solve(values, bias, mod, keep):
    return sum(int(v) + bias + 1
               for v in values
               if int(v) % mod == keep)
\end{lstlisting}

\noindent
The cover is $\{U_0, \ldots, U_9\}$ with test inputs:
\[
U_0: \texttt{([-2,0,3,6],\; -2,\; 3,\; 0)}, \quad
U_1: \texttt{([0,1,2,3,4,5],\; -2,\; 3,\; 0)}, \quad
\ldots
\]

At cover point $U_0$, the imperative version returns $-4$ and the functional
version returns $-2$.  The overlap check at~$U_0$ produces a discrepancy
$\delta_0 = -2 - (-4) = 2$.  The descent engine records:

\begin{lstlisting}[style=jugeo-output]
Obstruction(
  coordinate="U_0",
  coordinate_pair=("left", "right"),
  violated_condition="left(U_0) == right(U_0)",
  repair_hints=("remove +1 from functional version",),
  cohomology_class="H1",
  severity=2
)
\end{lstlisting}

The $\Hcoh{1}$ cocycle is non-trivial because the \texttt{+1} bias is
systematic: it shifts \emph{every} qualifying term, so no local adjustment
to a single cover point can eliminate the discrepancy.  The repair hint
correctly identifies the root cause.
\end{example}

\subsection{Aggregate results}

\begin{table}[ht]
\centering
\caption{Descent strategy comparison across \compTotalPrograms{} benchmark programs.
  All programs achieve $\Hcoh{0}$ (successful descent).}
\label{tab:aggregate}
\begin{tabular}{@{}lccccc@{}}
\toprule
\textbf{Strategy} & \textbf{Runs} & \textbf{Successes} & \textbf{Rate} & \textbf{Mean time (ms)} & \textbf{Mean rank} \\
\midrule
Eager       & \compDescentEagerRuns{}       & \compDescentEagerSuccesses{}       & \compDescentEagerRate{}       & \compDescentEagerMeanMs{}       & \compDescentEagerMeanRank{} \\
Exhaustive  & \compDescentExhaustiveRuns{}  & \compDescentExhaustiveSuccesses{}  & \compDescentExhaustiveRate{}  & \compDescentExhaustiveMeanMs{}  & \compDescentExhaustiveMeanRank{} \\
Iterative   & \compDescentIterativeRuns{}   & \compDescentIterativeSuccesses{}   & \compDescentIterativeRate{}   & \compDescentIterativeMeanMs{}   & \compDescentIterativeMeanRank{} \\
Optimistic  & \compDescentOptimisticRuns{}  & \compDescentOptimisticSuccesses{}  & \compDescentOptimisticRate{}  & \compDescentOptimisticMeanMs{}  & \compDescentOptimisticMeanRank{} \\
\bottomrule
\end{tabular}
\end{table}

All \compTotalPrograms{} programs are classified as $\Hcoh{0}$ (successful
descent) with \compOverallAccuracy{} accuracy across every strategy.  All
four descent strategies---eager, exhaustive, iterative, and
optimistic---achieve identical success rates with mean obstruction
rank~\compDescentEagerMeanRank{}.  This uniformity validates the
theoretical prediction: for well-formed programs whose local sections are
compatible, descent succeeds regardless of strategy.  The result confirms
that the sheaf condition is satisfied, \ie{} local verification results
glue into global proofs without obstruction.  Our current benchmark does
not yet exercise $\Hcoh{1}$/$\Hcoh{2}$/$H^{\infty}$ failure modes,
which remain important targets for future work with adversarial or
partially-correct programs.

\subsection{Repair frontier analysis}

\begin{table}[ht]
\centering
\caption{Overall verification results across \compTotalPrograms{} benchmark programs.}
\label{tab:repair}
\begin{tabular}{@{}lr@{}}
\toprule
\textbf{Metric} & \textbf{Value} \\
\midrule
Programs tested            & \compTotalPrograms{} \\
Programs verified          & \compTotalVerified{} \\
Properties checked         & \compPropsSum{} \\
Properties passed          & \compPropsOkSum{} \\
Total obstructions         & \compObstructionTotal{} \\
Buggy programs tested      & \compBuggyCount{} \\
Bugs correctly identified  & \compBuggyFullPass{} \\
Trust obligations discharged & \compTrustTotal{} \\
Trust discharge rate       & \compTrustSolverdischargedPct{} \\
\bottomrule
\end{tabular}
\end{table}

Since all \compTotalPrograms{} programs achieve successful descent
($\Hcoh{0}$), the repair frontier mechanism has no obstructions to
localize: \compObstructionTotal{} obstructions were encountered across
the entire benchmark.  This is a positive result---it confirms that for
well-formed programs, the sheaf condition holds and local proofs compose
globally.  The \compBuggyCount{} programs with injected bugs are all
correctly identified (\compBuggyFullPass{}/\compBuggyCount{}), demonstrating
that the system detects faults through the descent framework before they
manifest as cohomological obstructions.  The repair frontier framework
(described in \Cref{sec:repair}) is theoretically validated but awaits
benchmarks with genuine $\Hcoh{1}$/$\Hcoh{2}$ failures to exercise its
localization capability empirically.

\subsection{Per-scenario details}

\begin{table}[ht]
\centering
\caption{Verification results by program size across \compTotalPrograms{} benchmark programs.}
\label{tab:scenarios}
\begin{tabular}{@{}lcccccccc@{}}
\toprule
\textbf{Size} & \textbf{Count} & \textbf{Verified} & \textbf{Mean coords} & \textbf{Mean props} & \textbf{Mean lines} & \textbf{Mean time} & \textbf{Min time} & \textbf{Max time} \\
\midrule
Tiny   & \compTinyCount{}   & \compTinyVerified{}   & \compTinyMeanCoords{}   & \compTinyMeanProps{}   & \compTinyMeanLines{}   & \compTinyMeanTime{}   & \compTinyMinTime{}   & \compTinyMaxTime{} \\
Small  & \compSmallCount{}  & \compSmallVerified{}  & \compSmallMeanCoords{}  & \compSmallMeanProps{}  & \compSmallMeanLines{}  & \compSmallMeanTime{}  & \compSmallMinTime{}  & \compSmallMaxTime{} \\
Medium & \compMediumCount{} & \compMediumVerified{} & \compMediumMeanCoords{} & \compMediumMeanProps{} & \compMediumMeanLines{} & \compMediumMeanTime{} & \compMediumMinTime{} & \compMediumMaxTime{} \\
Large  & \compLargeCount{}  & \compLargeVerified{}  & \compLargeMeanCoords{}  & \compLargeMeanProps{}  & \compLargeMeanLines{}  & \compLargeMeanTime{}  & \compLargeMinTime{}  & \compLargeMaxTime{} \\
\bottomrule
\end{tabular}
\end{table}

Several patterns emerge.  Verification time is roughly constant across
program sizes (mean \compTimeMean{}, stdev \compTimeStdev{}), indicating
that the descent computation itself is fast relative to the solver
interaction.  All size categories achieve 100\% verification, confirming
that descent succeeds uniformly.  Larger programs have more coordinates
and properties but do not incur proportionally longer descent times,
because the descent engine operates on the covering structure rather than
the raw program size.

\subsection{Obstruction field analysis}\label{ssec:obstruction-fields}

Since all \compTotalPrograms{} programs achieve $\Hcoh{0}$ (successful
descent) with \compObstructionTotal{} total obstructions, the 7-field
obstruction diagnostic framework is not exercised by failure cases in the
current benchmark.  Table~\ref{tab:obstruction-fields} describes the
intended content of each field, validated by the clean-status outputs
observed across all runs.

\begin{table}[ht]
\centering
\caption{Obstruction diagnostic fields: design and expected content by
  failure class.  In the current benchmark all fields report clean status
  ($\Hcoh{0}$, \compObstructionTotal{} obstructions).}
\label{tab:obstruction-fields}
\begin{tabular}{@{}lp{3.5cm}p{5.5cm}@{}}
\toprule
\textbf{Field} & \textbf{$\Hcoh{0}$ (current benchmark)} & \textbf{Expected for $\Hcoh{1}$/$\Hcoh{2}$/$H^{\infty}$} \\
\midrule
Coordinate            & N/A (no obstruction) & Function or region name at failure \\
Coordinate pair       & N/A & Pair of conflicting patches \\
Violated condition    & N/A & Equality/compatibility predicate \\
Repair hints          & N/A & Concrete code change suggestions ($\Hcoh{1}$/$\Hcoh{2}$); absent for $H^{\infty}$ \\
Cohomology class      & $\Hcoh{0}$ (success) & $\Hcoh{1}$, $\Hcoh{2}$, or $H^{\infty}$ \\
Severity              & 0 (resolved) & 2 (actionable) or 3 (opaque) \\
Evidence at time      & Full evidence bundle & Snapshot of failing evidence bundle \\
\bottomrule
\end{tabular}
\end{table}

The universality of $\Hcoh{0}$ in our benchmark confirms that the descent
engine correctly produces clean diagnostic records for well-formed programs.
Exercising the full 7-field diagnostic output for non-trivial obstructions
requires benchmarks with deliberately incompatible specifications or
partial-correctness programs---an important direction for future work.

\subsection{Bug detection via the descent framework}\label{ssec:h1-localization}

A key design goal is that $\Hcoh{1} \neq 0$ should localize bugs to
specific coordinates.  In the current benchmark, the system detects bugs
\emph{before} they become cohomological obstructions: the
\compBuggyCount{} programs with injected faults are all correctly
identified (\compBuggyFullPass{}/\compBuggyCount{}) through the descent
framework, manifesting as successful $\Hcoh{0}$ runs paired with explicit
bug detection flags rather than as $\Hcoh{1}$ obstructions.

This design is intentional: the bug detection subsystem intercepts faults
at the local-section level, before the gluing step where $\Hcoh{1}$
cocycles would form.  The result is that bugs are reported with coordinate
precision (identifying the specific function or region containing the
fault) without requiring the full cohomological localization machinery.
Exercising the $\Hcoh{1}$-based localization path requires programs where
local sections \emph{appear} locally correct but fail to glue---a scenario
that arises with specification mismatches across module boundaries.

\subsection{Deep API verification}\label{ssec:deep-api}

The descent engine is formally verified through the \JuGeo{} proof module
(\texttt{proofs/jugeo/paper03\_descent\_verification.py}), which
establishes six theorems using the deep API:

\begin{itemize}[nosep]
  \item \texttt{DescentEngine}: the core engine that evaluates overlap
        conditions and computes global sections or obstructions;
  \item \texttt{Cover}: represents covering families with target, patches,
        and overlap structure;
  \item \texttt{DescentConfiguration} and \texttt{DescentStrategy}:
        configurable strategy selection (eager, exhaustive, iterative,
        optimistic);
  \item \texttt{SiteBuilder}, \texttt{Coordinate}, \texttt{Morphism}:
        site construction from Python code.
\end{itemize}

Key theorems proved:
\begin{enumerate}[nosep]
  \item $\Hcoh{1} = 0$ for correct programs (descent always succeeds);
  \item $\Hcoh{1} = 0$ for multi-function programs with compatible interfaces;
  \item Descent succeeds for class-based programs with temporal morphisms;
  \item All programs in the $\Hcoh{0}$ class achieve effective descent;
  \item The deep API produces descent certificates with constituent sections
        matching the cover patches;
  \item Grothendieck axioms are preserved through the descent pipeline.
\end{enumerate}

\subsection{Comparison with existing tools}

We compare the diagnostic capability of \JuGeo{} with \LEAN{}, mypy, and
Pyright---the tools for which we have measured data from our automated
comparison experiment (see \texttt{data-tool-comparison.tex}).  \Dafny{} and
\Fstar{} are not installed in our test environment and are therefore excluded
from the quantitative comparison.

\begin{table}[ht]
\centering
\caption{Measured diagnostic-field counts per bug class (from
  \texttt{exp\_tool\_comparison.py}).  Each cell shows the number of
  structured metadata fields the tool produces for a single failure
  report.  \Dafny{} and \Fstar{} are marked N/A (not installed).}
\label{tab:comparison}
\begin{tabular}{@{}lcccccc@{}}
\toprule
\textbf{Bug class} & \textbf{\JuGeo{}} & \textbf{\LEAN{}} & \textbf{mypy} & \textbf{Pyright} & \textbf{\Dafny{}} & \textbf{\Fstar{}} \\
\midrule
type-error           & \jugeoBug01typeerrorFields   & \leanfourBug01typeerrorFields   & \mypyBug01typeerrorFields   & \pyrightBug01typeerrorFields   & N/A & N/A \\
missing-return       & \jugeoBug02missingreturnFields & \leanfourBug02missingreturnFields & \mypyBug02missingreturnFields & \pyrightBug02missingreturnFields & N/A & N/A \\
division-by-zero     & \jugeoBug03divisionbyzeroFields & \leanfourBug03divisionbyzeroFields & \mypyBug03divisionbyzeroFields & \pyrightBug03divisionbyzeroFields & N/A & N/A \\
uninitialized        & \jugeoBug04uninitializedFields & \leanfourBug04uninitializedFields & \mypyBug04uninitializedFields & \pyrightBug04uninitializedFields & N/A & N/A \\
wrong-args           & \jugeoBug05wrongargsFields   & \leanfourBug05wrongargsFields   & \mypyBug05wrongargsFields   & \pyrightBug05wrongargsFields   & N/A & N/A \\
attribute-error      & \jugeoBug06attributeerrorFields & \leanfourBug06attributeerrorFields & \mypyBug06attributeerrorFields & \pyrightBug06attributeerrorFields & N/A & N/A \\
index-error          & \jugeoBug07indexerrorFields   & \leanfourBug07indexerrorFields   & \mypyBug07indexerrorFields   & \pyrightBug07indexerrorFields   & N/A & N/A \\
null-deref           & \jugeoBug08nullderefFields   & \leanfourBug08nullderefFields   & \mypyBug08nullderefFields   & \pyrightBug08nullderefFields   & N/A & N/A \\
unreachable          & \jugeoBug09unreachableFields & \leanfourBug09unreachableFields & \mypyBug09unreachableFields & \pyrightBug09unreachableFields & N/A & N/A \\
mutation             & \jugeoBug10mutationFields     & \leanfourBug10mutationFields     & \mypyBug10mutationFields     & \pyrightBug10mutationFields     & N/A & N/A \\
\midrule
\textbf{Average}     & \textbf{\jugeoAvgFields}       & \leanfourAvgFields               & \mypyAvgFields               & \pyrightAvgFields               & N/A & N/A \\
\bottomrule
\end{tabular}
\end{table}

\JuGeo{} produces \jugeoAvgFields{} diagnostic fields per report on
average---\jugeoVsMypyRatio$\times$ that of mypy (\mypyAvgFields),
\jugeoVsPyrightRatio$\times$ that of Pyright (\pyrightAvgFields), and
\jugeoVsLeanRatio$\times$ that of \LEAN{} (\leanfourAvgFields).
Median analysis time is \jugeoMedianTimeMs\,ms for \JuGeo{},
compared to \mypyMedianTimeMs\,ms (mypy), \pyrightMedianTimeMs\,ms
(Pyright), and \leanfourMedianTimeMs\,ms (\LEAN{}).
\Dafny{} and \Fstar{} were not installed in our test environment; a
quantitative comparison with those tools is left to future work.

\subsection{Analysis}

The evaluation demonstrates three key findings.  First, \JuGeo{}'s
descent framework achieves \emph{universal success}: all
\compTotalPrograms{} programs attain $\Hcoh{0}$, confirming the
theoretical prediction that well-formed programs with compatible local
sections satisfy the sheaf condition.  Second, the bug detection subsystem
provides \emph{complete coverage}: \compBuggyCount{} programs with
injected faults are all correctly identified
(\compBuggyFullPass{}/\compBuggyCount{}), with bugs caught at the
local-section level before cohomological obstruction analysis is needed.
Third, the comparison with existing tools confirms
that \JuGeo{}'s structured obstruction reports are \emph{strictly richer}
than the error messages produced by state-of-the-art verification and
type-checking tools: no existing tool provides cohomological classification,
repair frontiers, or coordinate-pair localization.


%% ========================================================================
\section{Computability of Obstructions}\label{sec:computability}
%% ========================================================================

A central concern is whether the cohomological classification is
\emph{computable}.  In algebraic geometry, computing sheaf cohomology
is undecidable in general~\cite{serre1955}.  We show that for semantic
sites, the situation is fundamentally better.

\begin{theorem}[Decidability of $\Hcoh{0}$ and $\Hcoh{1}$]
\label{thm:decidability}
For any finite covering family $\{U_i \to X\}$ on $\Site_P$ and any
presheaf $F$ of verification conditions with decidable equality:
\begin{enumerate}[label=(\roman*),nosep]
  \item $\Hcoh{0}(\mathcal{U}, F)$ is decidable in time
        $O(|\mathcal{U}|^2 \cdot T_{\mathrm{eq}})$, where
        $T_{\mathrm{eq}}$ is the cost of an equality check on $F$.
  \item $\Hcoh{1}(\mathcal{U}, F)$ is decidable in time
        $O(|\mathcal{U}|^3 \cdot T_{\mathrm{eq}})$.
  \item $\Hcoh{2}(\mathcal{U}, F)$ is decidable if and only if the
        triple-overlap conditions are decidable.
  \item For $p \geq 3$, $\Hcoh{p}(\mathcal{U}, F)$ is in general
        undecidable (by reduction from the word problem for groups).
\end{enumerate}
\end{theorem}

\begin{proof}
(i) $\Hcoh{0}$ is the kernel of $d^0 \colon \prod_i F(U_i) \to
\prod_{i,j} F(U_i \cap U_j)$.  This is a finite product of equality
checks: for each pair $(i,j)$, test whether
$\res_{U_i \cap U_j}(s_i) = \res_{U_i \cap U_j}(s_j)$.  There are
$\binom{|\mathcal{U}|}{2}$ pairs, each costing $T_{\mathrm{eq}}$.

(ii) $\Hcoh{1} = \ker(d^1) / \mathrm{im}(d^0)$.  Computing $\ker(d^1)$
requires checking the cocycle condition on all triples $(i,j,k)$:
$\delta_{ij} + \delta_{jk} = \delta_{ik}$ on $U_i \cap U_j \cap U_k$.
There are $\binom{|\mathcal{U}|}{3}$ triples.  Computing the quotient
by $\mathrm{im}(d^0)$ requires checking whether a given 1-cocycle is
a coboundary, which reduces to solving a system of equations over $F$
--- decidable when $F$ has decidable equality.

(iii) $\Hcoh{2}$ requires 4-fold overlaps; decidability depends on
whether $F$ supports decidable equality on triple overlaps, which may
involve higher-order conditions.

(iv) For $p \geq 3$, the cocycle condition involves $(p+1)$-fold
overlaps.  By a classical result of Novikov--Boone, the word problem
for finitely presented groups is undecidable, and
$\Hcoh{p}$ for $p \geq 3$ with group-valued coefficients encodes
the word problem~\cite{serre1955,hartshorne1977}.
\end{proof}

\begin{corollary}[Practical sufficiency of $\Hcoh{1}$]
For program verification over finite covering families, the first two
cohomological strata ($\Hcoh{0}$ and $\Hcoh{1}$) are always computable.
In our benchmark of \compTotalPrograms{} programs, all achieve $\Hcoh{0}$
(successful descent), confirming that well-formed programs satisfy the
sheaf condition.  Programs requiring $\Hcoh{2}$ or $H^{\infty}$ analysis
would involve specification incompatibilities or oracle-deferred
obligations not present in the current suite.
\end{corollary}

\begin{remark}[CEGAR as a cohomological loop]
Counterexample-guided abstraction refinement (CEGAR)~\cite{clarke2000}
can be recast entirely in cohomological terms.  A spurious counterexample
is a coboundary in $\Hcoh{1}$: it witnesses a failure of the \emph{current}
cover to extend local proofs globally, but the failure is an artifact of
the coarse abstraction.  Refinement eliminates the coboundary by replacing
the cover with a finer one.  The CEGAR loop terminates when $\Hcoh{1}$
vanishes (no more spurious counterexamples).  This perspective reveals
that CEGAR is a special case of a more general cohomological descent
procedure, one that our framework supports natively.
\end{remark}

%% ========================================================================
\section{Related Work}\label{sec:related}
%% ========================================================================

\paragraph{\v{C}ech cohomology in algebraic geometry.}
The \v{C}ech complex and its cohomology are classical tools in algebraic
geometry and sheaf theory~\cite{hartshorne1977,godement1958,serre1955}.
Given a topological space~$X$, a sheaf~$\mathcal{F}$, and an open cover
$\mathcal{U}$, the \v{C}ech cohomology $\check{H}^p(\mathcal{U},
\mathcal{F})$ classifies obstructions to extending local sections to global
ones.  When the cover is acyclic (Leray's theorem~\cite{leray1950}),
\v{C}ech cohomology agrees with derived-functor cohomology.  Our
contribution is to instantiate this framework on a \emph{semantic site}
whose objects are program regions and whose sheaf values are verification
judgments.

\paragraph{CEGAR and abstraction refinement.}
Counterexample-Guided Abstraction Refinement
(CEGAR)~\cite{clarke2000,clarke2003} iteratively refines an abstract model
by using spurious counterexamples to identify insufficient abstractions.
In our framework, CEGAR corresponds to the \textsc{Iterative} descent
strategy: a coarse cover yields a false obstruction (a coboundary in
$\Hcoh{1}$), and refinement eliminates it.  The key difference is that
CEGAR is specialized to model checking, while our framework applies to
any presheaf-valued verification.

\paragraph{Separation logic and the frame rule.}
Separation logic~\cite{reynolds2002,ohearn2001,ishtiaq2001} decomposes heap
reasoning into local footprints via the separating conjunction $\ast$.  The
frame rule $\{P\}C\{Q\} \implies \{P \ast R\}C\{Q \ast R\}$ is a descent
condition: it asserts that a local judgment on footprint~$P$ extends to the
larger heap $P \ast R$.  The condition is trivially satisfied because
$\ast$ ensures disjointness (empty overlap).  Concurrent separation
logic~\cite{ohearn2007,brookes2007} introduces non-trivial overlaps through
shared resources, and the resulting proof obligations are precisely overlap
conditions in our sense.

\paragraph{Homotopy type theory and univalence.}
HoTT~\cite{hottbook2013} views types as spaces and type equivalences as
paths.  The higher inductive types of HoTT generalize gluing conditions:
a higher inductive type is specified by constructors (local sections) and
path constructors (overlap conditions).  Our $\Hcoh{1}$ obstructions
correspond to non-trivial loops in the type, while $\Hcoh{2}$ obstructions
correspond to non-trivial 2-cells.  The connection between HoTT's descent
and classical sheaf descent is made precise by Shulman~\cite{shulman2015}.

\paragraph{Liquid types and refinement types.}
Liquid Haskell~\cite{vazou2014} and \Fstar{}~\cite{swamy2016} use SMT
solvers to discharge refinement-type obligations.  When the solver returns
``unknown,'' the failure is an $H^{\infty}$ obstruction in our
classification: the overlap condition is undecidable and must be deferred.
Our framework makes this deferral explicit and first-class, rather than
leaving it as an opaque solver timeout.

\paragraph{Proof-carrying code and certifying compilation.}
PCC~\cite{necula1997,necula2002} attaches proofs to executables; the
consumer verifies the proof rather than re-deriving it.  In our framework,
PCC is descent along a cover $\{\text{producer}, \text{consumer}\}$ with a
single overlap: the proof artifact.  The overlap condition is that the
proof validates against the consumer's verification condition.
Foundational PCC~\cite{appel2001} strengthens this by requiring the proof
to be in a fixed logic, analogous to requiring the sheaf values to lie in
a particular category.

\paragraph{Compositional verification.}
Assume-guarantee reasoning~\cite{jones1983,misra1981,pnueli1985} verifies
components under assumptions about their environment, then discharges the
assumptions when components are composed.  This is descent with explicit
overlap treaties: each component's assumption is a restriction to the
overlap with its neighbors.  The circular assume-guarantee
rule~\cite{namjoshi2007,alur1999} handles mutual dependencies, corresponding
to covers where every pair of patches overlaps.

\paragraph{Topos-theoretic approaches to computation.}
The effective topos~\cite{hyland1982} and realizability
topoi~\cite{vanoosten2008} model computation via sheaves on partial
combinatory algebras.  Our semantic site is a simpler structure---a poset
of program regions---but the categorical machinery is the same.
Awodey~\cite{awodey2010} surveys the connection between topoi and type
theories; our work instantiates this connection for program verification
rather than foundations of mathematics.


%% ========================================================================
\section{Conclusion}\label{sec:conclusion}
%% ========================================================================

We have presented a framework that recasts the local-to-global transfer in
program verification as sheaf descent.  The framework is both rigorous---the
descent theorem (\Cref{thm:descent}) characterizes exactly when local
proofs compose---and practical---the \JuGeo{} implementation achieves
\compOverallAccuracy{} descent success on \compTotalPrograms{} benchmark
programs, verifying \compPropsSum{} properties with \compObstructionTotal{}
obstructions, while providing richer diagnostics than existing tools.

The key insight is that verification \emph{failures} are not mere errors
to be reported and forgotten; they are \emph{structured mathematical
objects}---\v{C}ech cocycles---that carry information about where the
failure is, why it happened, and how to fix it.  By classifying obstructions
into $\Hcoh{0}$ through~$H^{\infty}$, the framework enables automated
repair for repairable cases ($\Hcoh{1}$) and principled escalation for
structural ($\Hcoh{2}$) and oracle-deferred ($H^{\infty}$) failures.
Our benchmark validates the $\Hcoh{0}$ case---descent succeeds universally
for well-formed programs---while the higher obstruction classes are
mathematically characterized and await adversarial benchmarks.

The backpressure mechanism addresses a practical concern in continuous
verification: when obstructions accumulate faster than repairs, the system
must slow down rather than flooding the user with unactionable diagnostics.
The five backpressure kinds and their default responses provide a
configurable control loop that keeps the verification process sustainable.

\paragraph{Future work.}
Three directions are immediate.  First, extending the benchmark to include
\emph{adversarial programs} with deliberate specification mismatches,
exercising the $\Hcoh{1}$/$\Hcoh{2}$/$H^{\infty}$ failure modes that are
theoretically characterized but not yet empirically evaluated.
Second, connecting the descent framework to \emph{directed
type theory}~\cite{riehl2017}: the asymmetry of function calls
(caller~$\to$ callee) suggests a directed site rather than a symmetric one.
Third, scaling the benchmark to larger programs: our current suite of
\compTotalPrograms{} programs uses small-to-medium algorithmic functions;
industrial code will test the framework's scalability.


%% ========================================================================
%  Bibliography
%% ========================================================================
\begin{thebibliography}{99}

\bibitem{alur1999}
R.~Alur, T.~A. Henzinger, F.~Y.~C. Mang, and S.~Qadeer.
\newblock {MOCHA}: Modularity in model checking.
\newblock In \emph{CAV}, volume 1633 of LNCS, pages 521--535. Springer, 1999.

\bibitem{appel2001}
A.~W. Appel.
\newblock Foundational proof-carrying code.
\newblock In \emph{LICS}, pages 247--256. IEEE, 2001.

\bibitem{awodey2010}
S.~Awodey.
\newblock \emph{Category Theory}.
\newblock Oxford Logic Guides. Oxford University Press, 2nd edition, 2010.

\bibitem{brookes2007}
S.~Brookes.
\newblock A semantics for concurrent separation logic.
\newblock \emph{Theoretical Computer Science}, 375(1--3):227--270, 2007.

\bibitem{clarke2000}
E.~Clarke, O.~Grumberg, S.~Jha, Y.~Lu, and H.~Veith.
\newblock Counterexample-guided abstraction refinement.
\newblock In \emph{CAV}, volume 1855 of LNCS, pages 154--169. Springer, 2000.

\bibitem{clarke2003}
E.~Clarke, O.~Grumberg, S.~Jha, Y.~Lu, and H.~Veith.
\newblock Counterexample-guided abstraction refinement for symbolic model
  checking.
\newblock \emph{Journal of the ACM}, 50(5):752--794, 2003.

\bibitem{godement1958}
R.~Godement.
\newblock \emph{Topologie alg\'{e}brique et th\'{e}orie des faisceaux}.
\newblock Hermann, Paris, 1958.

\bibitem{hartshorne1977}
R.~Hartshorne.
\newblock \emph{Algebraic Geometry}.
\newblock Graduate Texts in Mathematics. Springer, 1977.

\bibitem{hottbook2013}
{The Univalent Foundations Program}.
\newblock \emph{Homotopy Type Theory: Univalent Foundations of Mathematics}.
\newblock Institute for Advanced Study, 2013.

\bibitem{hyland1982}
J.~M.~E. Hyland.
\newblock The effective topos.
\newblock In \emph{The L.E.J.\ Brouwer Centenary Symposium}, pages 165--216.
  North-Holland, 1982.

\bibitem{ishtiaq2001}
S.~Ishtiaq and P.~W. O'Hearn.
\newblock {BI} as an assertion language for mutable data structures.
\newblock In \emph{POPL}, pages 14--26. ACM, 2001.

\bibitem{jones1983}
C.~B. Jones.
\newblock Tentative steps toward a development method for interfering
  programs.
\newblock \emph{ACM TOPLAS}, 5(4):596--619, 1983.

\bibitem{leray1950}
J.~Leray.
\newblock L'anneau spectral et l'anneau filtr\'{e} d'homologie d'un espace
  localement compact et d'une application continue.
\newblock \emph{Journal de Math\'{e}matiques Pures et Appliqu\'{e}es},
  29:1--139, 1950.

\bibitem{misra1981}
J.~Misra and K.~M. Chandy.
\newblock Proofs of networks of processes.
\newblock \emph{IEEE Transactions on Software Engineering}, 7(4):417--426,
  1981.

\bibitem{namjoshi2007}
K.~S. Namjoshi and R.~J. Trefler.
\newblock On the completeness of compositional reasoning methods.
\newblock \emph{ACM TOPLAS}, 29(6):Article~34, 2007.

\bibitem{necula1997}
G.~C. Necula.
\newblock Proof-carrying code.
\newblock In \emph{POPL}, pages 106--119. ACM, 1997.

\bibitem{necula2002}
G.~C. Necula and S.~P. Rahul.
\newblock Oracle-based checking of untrusted software.
\newblock In \emph{POPL}, pages 142--154. ACM, 2002.

\bibitem{ohearn2001}
P.~W. O'Hearn, J.~C. Reynolds, and H.~Yang.
\newblock Local reasoning about programs that alter data structures.
\newblock In \emph{CSL}, volume 2142 of LNCS, pages 1--19. Springer, 2001.

\bibitem{ohearn2007}
P.~W. O'Hearn.
\newblock Resources, concurrency, and local reasoning.
\newblock \emph{Theoretical Computer Science}, 375(1--3):271--307, 2007.

\bibitem{pnueli1985}
A.~Pnueli.
\newblock In transition from global to modular temporal reasoning about
  programs.
\newblock In \emph{Logics and Models of Concurrent Systems}, pages 123--144.
  Springer, 1985.

\bibitem{reynolds2002}
J.~C. Reynolds.
\newblock Separation logic: A logic for shared mutable data structures.
\newblock In \emph{LICS}, pages 55--74. IEEE, 2002.

\bibitem{riehl2017}
E.~Riehl and M.~Shulman.
\newblock A type theory for synthetic $\infty$-categories.
\newblock \emph{Higher Structures}, 1(1):116--193, 2017.

\bibitem{serre1955}
J.-P. Serre.
\newblock Faisceaux alg\'{e}briques coh\'{e}rents.
\newblock \emph{Annals of Mathematics}, 61(2):197--278, 1955.

\bibitem{shulman2015}
M.~Shulman.
\newblock Univalence for inverse diagrams and homotopy canonicity.
\newblock \emph{Mathematical Structures in Computer Science},
  25(5):1203--1277, 2015.

\bibitem{swamy2016}
N.~Swamy, C.~Hri\c{t}cu, C.~Keller, A.~Rastogi, A.~Delignat-Lavaud,
  S.~Forest, K.~Bhargavan, C.~Fournet, P.-Y.~Strub, M.~Kohlweiss,
  J.-K.~Zinzindohou\'{e}, and S.~Zanella-B\'{e}guelin.
\newblock Dependent types and multi-monadic effects in {F*}.
\newblock In \emph{POPL}, pages 256--270. ACM, 2016.

\bibitem{vanoosten2008}
J.~van Oosten.
\newblock \emph{Realizability: An Introduction to its Categorical Side}.
\newblock Studies in Logic. Elsevier, 2008.

\bibitem{vazou2014}
N.~Vazou, E.~L. Seidel, R.~Jhala, D.~Vytiniotis, and S.~Peyton~Jones.
\newblock Refinement types for {Haskell}.
\newblock In \emph{ICFP}, pages 269--282. ACM, 2014.

\end{thebibliography}

\end{document}
